\BourbakiExerciseGroup

\begin{enumerate}
\def\labelenumi{\arabic{enumi})}
\item
  设 E 是域 K 的一个扩张，x 与 y 是 K{[}X{]} 中同一不可约多项式在 E 中的两个不同的根。证明扩张 \(\mathbf{K}(x)\) 与 \(\mathbf{K}(y)\) 在 K 上不是线性不相交的（利用定理 2）。* 给出一个例子，其中 K = Q 且 \(\mathbf{K}(x) \cap \mathbf{K}(y) = \mathbf{K}\) （参见 V，第 161 页，习题 2）。*
\item
  设 \((\mathbf{E}_i)_{i\in I}\) 是域 \(\mathbf{K}\) 的一族扩张，都包含于 \(\mathbf{L}\) （ \(\mathbf{K}\) 的一个扩张）中。若 \(\mathbf{F}_i\) 是 \(\mathbf{K}\) 在 \(\mathbf{E}_i\) 中的相对代数闭包，证明 \(\mathbf{K}\) 在 \(\mathbf{E} = \bigcap_{i\in I}\mathbf{E}_i\) 中的相对代数闭包是 \(\mathbf{F} = \bigcap_{i\in I}\mathbf{F}_i\) .
\item
  证明域 K 的每个代数扩张 E 都与 \(K \times N\) 的一个子集等势（考虑把 E 的每个元素对应到它在 K 上的极小多项式的映射）。特别地，有限域的每个代数扩张都是可数的，而无限域 K 的每个代数扩张都与 K 等势。* 由此推出，在实数域 R 中存在素子域 Q 上的超越数，并且所有这些数组成的集合具有连续统的势。*
\item
  设 F 是域 E 的一个超越扩域。证明
\end{enumerate}

\[
[ \mathrm{F}: \mathrm{E} ] \geqslant \operatorname * {S u p} (\operatorname{Card} (\mathrm{E}), \operatorname{Card} (\mathbf {N}))
\]

若扩张 F 可由一族可数的元素生成，则等号成立。（利用 IV，第 89 页，习题 7。）

\begin{enumerate}
\def\labelenumi{\arabic{enumi})}
\setcounter{enumi}{4}
\tightlist
\item
  设 A 是一个交换代数，S 是 A 的一个子集，它作为 K-代数生成 A。假设 \(\operatorname{Card}(S) < \operatorname{Card}(K)\) 。证明：若 M 是 A 的一个极大理想，则域 A/M 是 K 的一个代数扩张。（当 \(\operatorname{Card}(K) \leqslant \operatorname{Card}(N)\) 时利用 V，第 174 页，习题 15，当
\end{enumerate}

\[
\operatorname{Card} (\mathbf {K}) > \operatorname{Card} (\mathbf {N}).
\]

时利用习题 4。）* 由此推出，\(\mathbf{A} / \mathfrak{M} = \mathbf{K}\) （当 \(\mathbf{K}\) 是代数闭域时）。*

* 6) 设 C 是一个不可数的代数闭域，且 \(\mathbf{P} = \mathbf{C}[(\mathbf{X}_i)_{i \in I}]\) 是 C 上的一个多项式代数。给定 P 的元素的两个族 \((f_\alpha)_{\alpha \in A}, (g_\beta)_{\beta \in B}\) ，假设：a) B 与 I 都是可数的。

\begin{enumerate}
\def\labelenumi{\alph{enumi})}
\setcounter{enumi}{1}
\tightlist
\item
  对于每个有限子集 \(\mathbf{A}'\) （ \(\mathbf{A}\) 的）与每个有限子集 \(\mathbf{B}'\) （ \(\mathbf{B}\) 的），存在 \((x_i)_{i\in I} \in C^I\) 使得
\end{enumerate}

\[
f _ {\alpha} (x _ {i}) = 0 \quad \text { and } \quad g _ {\beta} (x _ {i}) \neq 0 \quad \text { for   all } \quad \alpha \in \mathbf {A} ^ {\prime}, \beta \in \mathbf {B} ^ {\prime}.
\]

证明此时存在 \((x_{i})_{i\in I}\in \mathbf{C}^{I}\) 使得

\[
f _ {\alpha} (x _ {i}) = 0 \quad \text { and } \quad g _ {\beta} (x _ {i}) \neq 0 \quad \text { for   all } \quad \alpha \in A, \beta \in B.
\]

（设 \((\mathbf{Y}_{\beta})_{\beta \in \mathbb{B}}\) 是以 \(\mathbf{B}\) 为指标集的未定元族，并设 \(\mathbf{Q}\) 为多项式代数 \(\mathbf{C}[(\mathbf{X}_i)_{i\in I}, (\mathbf{Y}_{\beta})_{\beta \in \mathbb{B}}]\) 。设 \(\Lambda\) 为 \(\mathbf{Q}\) 关于由诸 \(f_{\alpha}\) 与诸 \(1 - g_{\alpha}\mathbf{Y}_{\alpha}\) 生成的理想所成的商。证明 \(\Lambda \neq 0\) （借助 \(a)\) 与 \(b)\) ），并对 \(\Lambda\) 的一个极大理想应用习题 5。）*

\begin{enumerate}
\def\labelenumi{\arabic{enumi})}
\setcounter{enumi}{6}
\tightlist
\item
  设 \(L\) 是一个域， \(A\) 是 \(L\) 的一个子环，且 \(K \subset L\) 是 \(A\) 的分式域.
\end{enumerate}

\begin{enumerate}
\def\labelenumi{\alph{enumi})}
\item
  证明：若 L 是有限生成的 A-模（II，第 206 页），则必有 A = K（考虑 K 在 L 中的补子空间（视 L 为 K 上的向量空间），证明 K 是有限生成的 A-模，从而必具有形式 \(s^{-1}A\) ，其中 \(s \in A\) ；最后证明 s 在 A 中必可逆）。
\item
  证明：若存在属于 L 的有限个元素 \(x_{j}\) ( \(1 \leqslant j \leqslant n\) )，它们在 K 上是代数的，使得 \(L = A[x_{1}, \ldots, x_{n}]\) ，则存在 A 中的元素 \(b \neq 0\) 使得 \(K = A[b^{-1}]\) （证明存在属于 A 的 \(b \neq 0\) ，使得 L 是有限生成的 \(A[b^{-1}]\) -模）。证明：要么 b 可逆（此时 K = A），要么 A 的每个非零理想都包含一个主理想 \(Ab^{m}\) 。反之亦然。域 k 上的形式幂级数环 \(k[[X]]\) 就是此类环的一个例子.
\end{enumerate}

* 8) 若 K 是域，证明 K 在域 K(X)（K 上单变元有理分式域）中是代数闭的（利用 K{[}X{]} 是主理想整环这一事实）。*

\begin{enumerate}
\def\labelenumi{\arabic{enumi})}
\setcounter{enumi}{8}
\tightlist
\item
  设 \(\mathbf{K}\) 是一个域， \(\mathbf{E}\) 是 \(\mathbf{K}\) 的一个扩张，且 \(\mathbf{S}\) 是 \(\mathbf{E} - \mathbf{K}\) 的一个子集.
\end{enumerate}

\begin{enumerate}
\def\labelenumi{\alph{enumi})}
\item
  证明由扩张 \(\mathbf{L} \subset \mathbf{E}\) （ \(\mathbf{K}\) 的、满足 \(\mathbf{L} \cap \mathbf{S} = \varnothing\) 的）组成的集合按包含关系排序具有一个极大元 \(\mathbf{M}\) .
\item
  证明：若 S 是有限的，则 E 是 M 的代数扩张。
\end{enumerate}

* 10) 设 K 是一个域，P 是环 K{[}X, Y{]} 中两个未定元的多项式。假设存在多项式 H ∈ K{[}X{]} 以及两个多项式 Q, R ∈ K{[}X, Y{]} 使得 H(X) P(X, Y) = Q(X, Y) R(X, Y)，并且 Q(X, Y) 的系数（作为 Y 的多项式）是 K{[}X{]} 中的多项式，且没有非常数的公因子。证明：R(X, Y) 的所有系数（作为 Y 的多项式）都能被 H(X) 整除。（将 H, P, Q, R 视为域 K(Y) 上关于 X 的多项式，并将它们在 Ω{[}X{]} 中分解为一次因式，其中 Ω 是 K(Y) 的一个代数闭包；注意，若 H(X) 与 Q(X, Y) 有一个一次公因子，则该因子将整除 Q(X, Y) 的所有系数（作为 Y 的多项式），并利用事实（IV，第 12 页）：若 K' 是 K 的一个扩域，则 K'{[}X{]} 中两个 K{[}X{]} 的多项式的最大公因式属于 K{[}X{]}。）*

\begin{enumerate}
\def\labelenumi{\arabic{enumi})}
\setcounter{enumi}{10}
\tightlist
\item
  设 \(\mathbf{E}\) 是域 \(\mathbf{K}, x \in \mathbf{E}\) 是 \(\mathbf{K}\) 上的一个超越元，从而 \(\mathbf{K}(x)\) 同构于 \(\mathbf{K}(\mathbf{T})\) （即 \(\mathbf{K}\) 上关于一个未定元的有理分式域）。于是每个 \(y \in \mathbf{K}(x)\) 都可写成 \(y = g(x) / h(x)\) ，其中 \(g\) 与 \(h\) 是 \(\mathbf{K}[\mathbf{X}]\) 的两个互素的多项式（IV，第 12 页），在相差 \(\mathbf{K}\) 中一个非零公因子的意义下确定；\(g\) 与 \(h\) 的次数中的较大者称为 \(y\) 关于 \(x\) 的高度.
\end{enumerate}

\begin{enumerate}
\def\labelenumi{\alph{enumi})}
\tightlist
\item
  证明在多项式环 \(\mathbf{K}(y)[\mathbf{X}]\) 中，多项式 \(g(\mathbf{X}) - yh(\mathbf{X})\) 当 \(y \notin \mathbf{K}\) 时是不可约的。由此推出，若 \(y\) 的高度为 \(n > 0\) （关于 \(x\) 的），则 \(\mathbf{K}(x)\) 是一个 \(n\) 次的、关于 \(\mathbf{K}(y)\) 的代数扩张.
\end{enumerate}

设 \(\mathrm{P}(\mathbf{X}) = \sum_{i=0}^{n} y_i \mathbf{X}^i\) 为 \(x\) 在 \(\mathrm{K}(y)\) 上的极小多项式。证明对于

\(0 \leqslant i \leqslant n\) ，我们或者有 \(y_i \in \mathbf{K}\) ，或者有 \(\mathbf{K}(y_i) = \mathbf{K}(y)\) .

\begin{enumerate}
\def\labelenumi{\alph{enumi})}
\setcounter{enumi}{1}
\item
  由 a) 推出，每个 \(y \in \mathbf{K}(x)\) 若满足 \(\mathbf{K}(y) = \mathbf{K}(x)\) ，则具有形式 \((ax + b)/(cx + d)\) ，其中 a、b、c、d 是 K 中满足 \(ad - bc \neq 0\) 的元素；反之亦然。求出 \(\mathbf{K}(x)\) 的全部 K-自同构.
\item
  证明：若 \(y \in \mathbf{K}(x)\) 的高度为 \(n\) （关于 \(x\) 的），且 \(z \in \mathbf{K}(y)\) 的高度为 \(m\) （关于 \(y\) 的），则 \(z\) 的高度为 \(mn\) （关于 \(x\) 的）.
\end{enumerate}

* \(d\) ) 设 F 是 K 的一个扩域，使得 \(\mathbf{K} \subset \mathbf{F} \subset \mathbf{K}(x)\) 且 \(\mathbf{F} \neq \mathbf{K}\) ；设 y 是 F 中的一个元素，其高度 \(m\) （关于 \(x\) ）取到可能的最小值；证明 \(\mathbf{F} = \mathbf{K}(y)\) （“吕罗特定理”）。（设 P 为 \(x\) 在 F{[}T{]} 中的极小多项式；证明 \(\mathrm{P(T)} = \mathrm{Q(T,x)}/\mathrm{R}(x)\) ，其中 Q 是 K{[}T,X{]} 中的一个多项式，它关于 X 的次数 \(\geqslant m\) ，且不被 K{[}X{]} 的任何非常数多项式整除，而 R 是 K{[}X{]} 的一个多项式。若 \(y = g(x)/h(x)\) ，其中 \(g\) 与 \(h\) 在 K{[}X{]} 中互素，注意由 \(a\) )，多项式 \(g(T)h(X) - g(X)h(T)\) 不被 K{[}T{]} 或 K{[}X{]} 的任何非常数多项式整除；由此必然推出

\[
g (T) h (X) - g (X) h (T) = c. Q (T, X), \quad \text { where } c \in K,
\]

利用习题10。）*

* 12) a) 采用习题 11 中的记号与假设，证明：K 的扩张 F 满足 \(\mathbf{K} \subset \mathbf{F} \subset \mathbf{K}(x)\) ，要使 \(\mathbf{F} \cap \mathbf{K}[x] \neq \mathbf{K}\) ，其充要条件是 \(\mathbf{F} = \mathbf{K}(y)\) ，其中 \(y = g(x)\) ，\(g \in \mathbf{K}[\mathrm{T}]\) 为非常数多项式；此时我们有 \(\mathbf{K}(y) \cap \mathbf{K}[x] = \mathbf{K}[y]\) 。 （写 \(\mathbf{F} = \mathbf{K}(y)\) ，其中 \(y = g(x)/h(x)\) , \(g\) 与 \(h\) 是 \(\mathbf{K}[\mathrm{T}]\) 的互素多项式。如果存在一个非常数多项式 \(\mathbf{P} \in \mathbf{K}[\mathrm{T}]\) 使得 \(\mathbf{P}(x) \in \mathbf{F}\) ，我们可以写出 \(\mathbf{P}(x) = \mathbf{Q}(y)/\mathbf{R}(y)\) ，其中 \(\mathbf{Q}\) 与 \(\mathbf{R}\) 是 \(\mathbf{K}[\mathrm{T}]\) 的两个互素的首一多项式，并非两者皆为常数。分两种情况讨论，视 \(\mathbf{R} = 1\) 或 \(\mathbf{R}\) 为非常数而定；在第二种情形下，将 \(\mathbf{Q}\) 与 \(\mathbf{R}\) 在 \(\Omega\) （ \(\mathbf{K}\) 的一个代数闭包）中分解为一次因式，并注意到，如果 \(a, b\) 是 \(\Omega\) 的不同元素，则 \(g - ah\) 与 \(g - bh\) 在 \(\Omega[\mathrm{T}]\) 中互素。）

\begin{enumerate}
\def\labelenumi{\alph{enumi})}
\setcounter{enumi}{1}
\item
  如果 \(\mathbf{F} = \mathbf{K}(y)\) ，其中 \(y = g(x)\) ，这里 \(g \in \mathbf{K}[T]\) 是非常数的，则我们可以假设 \(g(T) = Tg_1(T)\) ，其中 \(g_1 \in \mathbf{K}[T]\) 。证明如果 \(\mathbf{K}(xg(x)) \cap \mathbf{K}[y] \neq \mathbf{K}\) ，我们必然有 \(g_1(T) = aT^n\) ，其中 \(a \in \mathbf{K}\) 且 \(n > 0\) 。（注意，根据假设及 \(a\) ），存在两个非常数多项式 \(P, Q\) （在 \(K[T]\) 中）使得 \(P(xg(x)) = Q(g(x)) \notin K\) ；由此推出存在两个元素 \(a, b\) （在 \(K\) 中）以及一个整数 \(m\) 使得 \(g(T)\) 整除 \(a - bT^m\) 。） c) 设 \(y, z\) 为属于 \(K[x]\) 的两个元素，使得 \(K(y) \cap K(z) \neq K\) ；证明此时 \(K[y] \cap K[z] \neq K\) .
\item
  求 \(y \in \mathbf{K}(x)\) 使得 \(\mathbf{K}(y) = \mathbf{K}(x)\) 但是 \(\mathbf{K}[y] \cap \mathbf{K}[x] = \mathbf{K}\) . *
\end{enumerate}

* 13) 设 K 为一个域， \(P(T, X) = a_0(T)X^n + a_1(T)X^{n-1} + \cdots + a_n(T)\) 为 K{[}T, X{]} 中的一个多项式；假设： \(1^\circ\) 多项式 \(a_0\) 不被 T 整除； \(2^\circ\) 多项式 \(a_1, \ldots, a_n\) 中的每一个都被 T 整除； \(3^\circ\) 多项式 \(a_n\) 不被 \(T^2\) 整除。证明在这些条件下，P 作为 K(T){[}X{]} 中的多项式是不可约的（用反证法，利用习题 10）。*

* 14：采用习题 11 中的记号，证明如果我们取 \(g(X) = X\) 且 \(h(X) = X^n + X + 1 \quad (n \geqslant 2)\) ，则多项式 \((g(X)h(x) - h(X)g(x)) / (X - x)\) （作为 \(K(x)[X]\) 的元素）是不可约的（利用习题 13）。*

* 15）给出一个域 K 的例子，它包含两个子域 \(\mathbf{K}_1\) , \(\mathbf{K}_2\) ，使得 K 在 \(\mathbf{K}_1\) 上是代数的，在 \(\mathbf{K}_2\) 上也是代数的，但在 \(\mathbf{K}_1 \cap \mathbf{K}_2\) 上不是代数的。 （取 \(\mathbf{K} = \mathbf{Q}(T)\) , \(\mathbf{K}_1 = \mathbf{Q}(T^2)\) ，以及 \(\mathbf{K}_2 = \mathbf{Q}(T^2 - T)\) 。验证 \([\mathbf{K} : \mathbf{K}_1] = 2\) , \([\mathbf{K} : \mathbf{K}_2] = 2\) 且 \(\mathbf{K}_1 \cap \mathbf{K}_2 = \mathbf{Q}\) 。 为此，注意一个元素 \(f(T)\) （ \(\mathbf{K}_1 \cap \mathbf{K}_2\) 中的）满足条件 \(f(-T) = f(T)\) 与 \(f(1 - T) = f(T)\) .)

\BourbakiExerciseGroup

\begin{enumerate}
\def\labelenumi{\arabic{enumi})}
\tightlist
\item
  设 K 为一个域，且 \((\mathrm{E}_{i})_{i\in I}\) 为 K 的任意一族扩张。证明存在 K 的一个扩张 E，并且对每个 \(i\in I\) ，存在一个 K-同构 \(u_{i}\) （从 \(E_{i}\) 到 E 中的），使得 E 由诸 \(u_{i}(\mathrm{E}_{i})\) 的并生成（如命题 4（V，第 22 页）中那样论证）。
\end{enumerate}

* 2) 设 \(\mathbf{K}\) 是特征为 0 的代数闭域， \(\mathrm{F} = \mathrm{K}((\mathbf{X}^{1 / n!}))\) 为 \(\mathbf{X}^{1 / n!}\) 的形式幂级数域，而 \(\mathrm{E} = \mathrm{K}((\mathbf{X}))\) ，其中 \(\mathbf{X} = (\mathbf{X}^{1 / n!})^{n!}\) ，是 \(\mathbf{F}\) 中由 \(\mathbf{X}\) 的形式幂级数组成的子域。对于每个形式幂级数 \(f \in \mathbf{K}((\mathbf{X}))\) ，我们用 \(\omega(f)\) 表示它的阶。

\begin{enumerate}
\def\labelenumi{\alph{enumi})}
\item
  设 \(P(Y) = a_{0}Y^{n} + a_{1}Y^{n-1} + \cdots + a_{n}\) 是 E{[}Y{]} 中次数为 n 的多项式，使得 \(a_{n} \neq 0\) 。证明存在整数的一个严格递增序列 \((i_{k})_{0 \leqslant k \leqslant r}\) （在区间 \([0, n]\) 中），使得： \(1^{\circ} i_{0} = 0\) , \(i_{r} = n\) ; \(2^{\circ} \omega(a_{i_{k}})\) 对 \(0 \leqslant k \leqslant r\) 是有限的; \(3^{\circ}\) 对于每个满足 \(0 \leqslant j \leqslant n\) 、异于诸 \(i_{k}\) 且使 \(\omega(a_{j})\) 为有限的指标 j，点 \((j, \omega(a_{j})) \in \mathbf{R}^{2}\) 位于经过点 \((i_{k}, \omega(a_{i_{k}}))\) , \((i_{k-1}, \omega(a_{i_{k-1}}))\) 的直线上方，且当 \(j < i_{k-1}\) 或 \(j > i_{k} (1 \leqslant k \leqslant r)\) 时严格位于其上方。连接点 \((i_{k-1}, \omega(a_{i_{k-1}}))\) 与 \((i_{k}, \omega(a_{i_{k}}))\) （对 \(1 \leqslant k \leqslant r\) ）的诸线段之并称为 P 的牛顿多边形，前述线段称为其边，而点 \((i_{k}, \omega(a_{i_{k}}))\) 为其顶点。
\item
  令 \(\rho_{k} = i_{k} - i_{k - 1}\) , \(\sigma_{k} = (\omega (a_{i_{k}}) - \omega (a_{i_{k - 1}})) / \rho_{k}\) ；证明 \(P\) 恰有 \(\rho_{k}\) 个零点（按正确的重数计算）属于域 \(K((X^{1 / (\rho_k!)}))\) 且其阶等于 \(\sigma_{k}\) （用递归方法确定满足方程 \(P(z) = 0\) 的这种类型的形式幂级数的系数）.
\item
  得出结论：诸域 \(\mathbf{K}((\mathbf{X}^{1 / n}))\) （对所有整数 \(n > 0\) ）的并集是域 \(\mathrm{E} = \mathbf{K}((\mathbf{X}))\) 的一个代数闭包（“皮瑟定理”）。
\item
  如果 \(\mathbf{K}'\) 是特征为 2 的域，并且 \(\mathbf{E}' = \mathbf{K}'((\mathbf{X}))\) ，则多项式 \(\mathbf{Y}^2 + \mathbf{X}\mathbf{Y} + \mathbf{X}\) （ \(\mathbf{E}'[\mathbf{Y}]\) 中的）在诸域 \(\mathbf{K}'((\mathbf{X}^{1/n!}))\) 的并集中没有根。 *
\end{enumerate}

\BourbakiExerciseGroup

\begin{enumerate}
\def\labelenumi{\arabic{enumi})}
\tightlist
\item
  \begin{enumerate}
  \def\labelenumii{\alph{enumii})}
  \tightlist
  \item
    设E是域K的一个扩域， \(x \in E\) 是K上的一个超越元，并设 \(F = K(x^{n})\) ，其中 n 为整数，n \textgreater{} 1. 证明：如果交集 \(F \cap K[(x - 1)(x^{n} - 1)]\) 不同于K，则K的特征为p \textgreater{} 0，并且 \(n = p^{e}\) ，从而 \(K(x)\) 是F的一个p-根式扩张。（令y = x - 1，并利用V，第149页，习题12，b），证明 \((y + 1)^{n} - 1 = ay^{n}\) 对某个 \(a \in K\) 成立.)
  \end{enumerate}
\end{enumerate}

\begin{enumerate}
\def\labelenumi{\alph{enumi})}
\setcounter{enumi}{1}
\item
  证明如果 \(\mathbf{K} \subset \mathbf{F}' \subset \mathbf{K}(x)\) ，如果 \(\mathbf{K}(x)\) 不是 \(p\) -根式扩张（在 \(\mathbf{F}'\) 上），并且如果 \(\mathbf{K} \neq \mathbf{F}' \neq \mathbf{K}(x)\) ，则存在 \(z \in \mathbf{K}[x]\) 使得 \(z \notin \mathbf{K}\) 且 \(\mathbf{F}' \cap \mathbf{K}[z] = \mathbf{K}\) 。（用反证法，利用a）以及V，第 149 页，习题 12，b）。）
\item
  在b)中相同的假设下，假设 \(F' = K(y)\) ，其中 \(y \in K[x]\) ；证明存在 \(u \in K[x]\) 使得 \(u \notin K\) 且 \(F' \cap K(u) = F' \cap K[u] = K\) （注意，若此命题不成立，则应有 \(F' \cap K[z] \neq K\) 对所有 \(z \in K[x]\) 满足 \(z \notin K\) ，利用V，第 149 页，习题 12，c）。）
\end{enumerate}

\begin{enumerate}
\def\labelenumi{\arabic{enumi})}
\setcounter{enumi}{1}
\item
  设 K 是特征为 p 的域，p \textgreater{} 0，且 f 是 K{[}X{]} 中的一个不可约的首一多项式。证明在 K{[}X{]} 中，多项式 \(f(\mathbf{X}^{p})\) 要么不可约，要么是某个不可约多项式的 p 次幂，这取决于 f 是否存在一个不属于 \(K^{p}\) 的系数（把 \(f(\mathbf{X}^{p})\) 在 \(\Omega[X]\) 中分解为一次因式，其中 \(\Omega\) 是 K 的一个代数闭扩张）。
\item
  设 K 是特征为 p 的域，p \textgreater{} 2，设F为有理分式域 \(\mathrm{K}(\mathbf{X}, \mathbf{Y})\) 。设 \(\mathrm{E} = \mathrm{F}(\theta)\) 为 F 的由以下多项式的一个根 \(\theta\) 所生成的扩张
\end{enumerate}

\[
f (\mathbf {Z}) = \mathbf {Z} ^ {2 p} + \mathbf {X Z} ^ {p} + \mathbf {Y}
\]

（它是 F{[}Z{]} 中的多项式）。证明 \([E:F(\theta^{p})]=p\) ，但 E 不包含不属于 F 的 F 上的 p-根式元。（首先注意 f 在 F{[}Z{]} 中不可约；如果存在 \(\beta\in E\) 使得 \(\beta^{p}\in F\) , \(\beta\notin F\) ，则 f 在 F( \(\beta\) ){[}Z{]} 中将可约；利用习题2，证明此时 \(X^{1/p}\) 与 \(Y^{1/p}\) 将属于 E，并且我们将有 \([E:F]\geqslant p^{2}\) 。)（参见第V章第7节习题1。）

\BourbakiExerciseGroup

\begin{enumerate}
\def\labelenumi{\arabic{enumi})}
\item
  设 K 是一个域，A 是一个交换 K-代数。要使 A 是平展的，其充要条件是：存在 A 的一个有限子集 S，它生成 K-代数 A，并且使得对所有 \(x \in S\) ，由 x 生成的子代数 \(K[x]\) 是平展的。（若 S 生成 A，则 K-代数 A 同构于诸 \(K[x]\) （其中 x 遍历 S）的张量积的一个商。）
\item
  设 \(\Gamma\) 是一个有限交换幺半群。考虑以下条件 \(E_0\) 与 \(E_p\) ，其中 \(p\) 表示一个素数。
\end{enumerate}

(E0) 如果 \(x \in \Gamma\) 并且如果 \(e, d\) 是整数 \(\geqslant 0\) ，则关系 \(x^{e + d} = x^d\) 蕴含关系 \(x^{e + 1} = x\) .

\((E_p)\) 如果 \(x, y \in \Gamma\) ，则关系 \(x^p = y^p\) 蕴含 \(x = y\) .

设 \(\mathbf{K}\) 是特征为 \(p (\geqslant 0)\) 的域，并令 \(\mathbf{A}\) 为 K-代数 \(\mathbf{K}^{(\Gamma)}\) （幺半群 \(\Gamma\) 的）. \(a)\) 要使 \(\mathbf{A}\) 是平展的，其充分必要条件为条件 \((\mathrm{E}_p)\) 成立。（充分性利用习题 1。对于必要性，当 \(p = 0\) 时，注意 \(x^{e+1} - x\) 在 \(x^{e+d} = x^d\) 成立时是幂零的；对于 \(p > 0\) 利用 V，第 35 页的推论。）

\begin{enumerate}
\def\labelenumi{\alph{enumi})}
\setcounter{enumi}{1}
\item
  设 X 为 \(\Gamma\) 到 K 的乘法幺半群中的同态组成的集合。我们有 Card \(X \leqslant\) Card \(\Gamma\) 。若 K 是代数闭的，则 Card X = Card \(\Gamma\) 成立当且仅当条件 \((E_{p})\) 被满足。
\item
  如果 \(\Gamma\) 是一个群，条件 \((\mathbf{E}_0)\) 总是成立；如果 \(p > 0\) ，条件 \((\mathbf{E}_p)\) 等价于说 \(\Gamma\) 的每个元素的阶与 \(p\) 互素，或者等价地说 Card \(\Gamma\) 不被 \(p\) 整除。
\end{enumerate}

\BourbakiExerciseGroup

\begin{enumerate}
\def\labelenumi{\arabic{enumi})}
\tightlist
\item
  \begin{enumerate}
  \def\labelenumii{\alph{enumii})}
  \tightlist
  \item
    设 \(\mathbf{K}\) 是特征为 \(p > 0\) 的域。一个有限次的代数扩张 \(\mathbf{F}\) （\(\mathbf{K}\) 上的）被称为例外的，如果它在 \(\mathbf{K}\) 上不是可分的，并且它不包含任何元素 \(x \notin K\) 是 \(p\) -根式的（在 \(K\) 上），换句话说，如果 \(F^p \cap K = K^p\) （V，第151页，习题3）。证明如果 \(F\) 是例外的并且如果 \(E = F_s\) , \(z \in F\) , \(z \notin E\) 且 \(z^p \in E\) ，则对每个 \(t \in K\) 我们有 \(z^p \notin E^p(t)\) （注意，否则我们会有 \(E^p(t) = E^p(t^p)\) 因此 \(t \in F^p\) ).
  \end{enumerate}
\end{enumerate}

\begin{enumerate}
\def\labelenumi{\alph{enumi})}
\setcounter{enumi}{1}
\tightlist
\item
  反之，设E是K的有限次可分代数扩张，并设 \(x \in E\) 使得对于每个 \(t \in K\) 我们有 \(x \notin E^p(t)\) 。设 \(F = E(x^{1/p})\) ，从而 \(E = F_s\) 。证明 F 是 K 的例外扩张。（如果存在 \(y \in F\) 使得 \(y \notin E\) 且 \(y^p = z \in K\) ，证明我们将有 \(z \in E^p(x)\) ；这将蕴含 \(z \in E^p\) ，因此 \(y \in E\) ，这是一个矛盾。）
\end{enumerate}

K 的有限次可分扩张 E 称为强可分的，如果域 \(E^{p}(t)\) （其中 t 遍历 K）的并集不同于 E；因此这等同于说 \(E = F_{s}\) ，其中 F 是 K 的一个不可分例外扩张。K 的强可分扩张不可能是完全域。

\begin{enumerate}
\def\labelenumi{\alph{enumi})}
\setcounter{enumi}{2}
\tightlist
\item
  证明如果 E 是 K 的强可分扩张，则不可能 \([E:E^{p}]=p\) （用反证法论证，表明 \([E:E^{p}]=p\) 将意味着 \(\mathrm{E}^{p}(t)=\mathrm{E}\) 对所有 \(t\in K\) 使得 \(t\notin K^{p}\) ；另一方面注意，如果 \(K\subset E^{p}\) ，则除非 E 是完全的，否则 E 在 K 上不能是强可分的。)
\end{enumerate}

\(d\) ) 证明如果 \([\mathbf{K}:\mathbf{K}^p ] = p\) ，则不存在 \(\mathbf{K}\) 的强可分扩张，因此也不存在 \(\mathbf{K}\) 的例外扩张。

\begin{enumerate}
\def\labelenumi{\arabic{enumi})}
\setcounter{enumi}{1}
\tightlist
\item
  设 \(k\) 是特征为 \(p, a \in k\) 的非完全域，其中 \(a \notin k^p\) ， \(F\) 是有理分式域 \(k(X)\) （ \(k\) 上关于一个未定元的）；令 \(y = X^{p^2} / (X^p + a)\) 与 \(K = k(y)\) .
\end{enumerate}

\begin{enumerate}
\def\labelenumi{\alph{enumi})}
\item
  证明： \([\mathbf{F}:\mathbf{K}] = p^2\) ，从而 \(\mathrm{T}^{p^2} - y\mathrm{T}^p - ay\) 是 \(\mathbf{X}\) 在 \(\mathbf{K}[\mathrm{T}]\) 中的极小多项式；由此推出 \(\mathbf{K}\) 在 \(\mathbf{F}\) 中的相对可分闭包是 \(\mathrm{E} = k(\mathbf{X}^p)\) （参见 V，第 148 页，习题 11）。
\item
  证明 F 是 K 的例外扩张（习题 1）。（用反证法，假设 \(b \in F\) 使得 \(b \notin K\) 且 \(b^{p} \in K\) 。注意 \(\mathbf{K}(b) = k(z)\) ，其中 \(z = r(\mathbf{X}) / s(\mathbf{X})\) ，满足 \(\deg r = p\) , \(\deg s \leqslant p\) 且 r、s 是 \(k[X]\) 中的互素多项式；证明 \(z^{p} \in K\) 并且X在K上的极小多项式为 \(r(T)^{p} - z^{p}s(T)^{p}\) ，相差 K 中的一个因子。利用a)得出结论，我们将有 \(a \in k \cap F^{p} = k^{p}\) ，与假设相反。）
\end{enumerate}

\begin{enumerate}
\def\labelenumi{\arabic{enumi})}
\setcounter{enumi}{2}
\tightlist
\item
  设 K 是特征为 p 的域，p \textgreater{} 0，E是K的有限次扩张， \(E_{r}\) 与 \(E_{s}\) 是包含在 E 中的 K 的最大 p-根式扩张与可分扩张。证明 \(E_{s} \otimes_{K} E_{r}\) 同构于 \(E_{r}\) 在E中的相对可分闭包，并与之等同；若它不等于E，则E是 \(E_{r}\) 的例外扩张，且 \(E_{s} \otimes_{K} E_{r}\) 是 \(E_{r}\) 的强可分扩张。
\end{enumerate}

反之，设E是域L的一个例外扩张，并设K是L的一个子域，使得L在K上是p-根式的，且 \([L:K]<+\infty\) ；则 L 是包含于 E 中的 K 的最大 p-根式扩张。

\begin{enumerate}
\def\labelenumi{\arabic{enumi})}
\setcounter{enumi}{3}
\tightlist
\item
  设 \(\mathbf{K}\) 是一个域， \(\mathbf{E} = \mathbf{K}(x)\) 是 \(\mathbf{K}\) 的一个可分代数扩张，且 \(f\) 是 \(x\) 在 \(\mathbf{K}[\mathbf{X}]\) 中的极小多项式.
\end{enumerate}

\begin{enumerate}
\def\labelenumi{\alph{enumi})}
\item
  设 \(\mathbf{F}\) 是 \(\mathbf{K}\) 的一个扩张；在环 \(\mathbf{F}[\mathbf{X}]\) 中，多项式 \(f\) 分解为乘积 \(f_{1}f_{2}\ldots f_{r}\) ，各因子是不可约、可分且两两互异的多项式。证明代数 \(\mathbf{E} \otimes_{\mathbf{K}} \mathbf{F}\) 同构于域的乘积 \(\mathbf{F}[\mathbf{X}] / (f_j)\) ( \(1 \leqslant j \leqslant r\) ).
\item
  假设 \(\mathbf{F} = \mathbf{K}(y)\) 是 \(\mathbf{K}\) 的可分代数扩张；设 \(g\) 为 \(y\) 在 \(\mathbf{K}[\mathbf{X}]\) 中的极小多项式，并设 \(g = g_1g_2\ldots g_s\) 为 \(g\) 在E{[}X{]}中分解成不可约多项式的乘积。证明 r = s，并且如果 \(m = \deg f\) , \(n = \deg g\) , \(m_j = \deg f_j\) , \(n_j = \deg g_j\) ，则我们在必要时对 \(g_j\) 作置换后，可假设 \(m_j / n_j = m / n\) 对 \(1 \leq j \leq r\) 成立.
\end{enumerate}

\begin{enumerate}
\def\labelenumi{\arabic{enumi})}
\setcounter{enumi}{4}
\tightlist
\item
  设 \(\mathbf{K}\) 是有 \(q\) 个元素的有限域，并令 \(\mathbf{A}\) 为 K-代数 \(\mathbf{K}^n\) 。要使 \(\mathbf{A}\) 由单个元素生成，必要且充分条件是 \(n \leqslant q\) .
\end{enumerate}

\BourbakiExerciseGroup

\begin{enumerate}
\def\labelenumi{\arabic{enumi})}
\item
  设 \(\mathbf{K}\) 是一个交换环， \(f(\mathbf{X})\) 是 \(\mathbf{K}[\mathbf{X}]\) 中的一个首一多项式， \(\mathbf{A}\) 是 K-代数 \(\mathbf{K}[\mathbf{X}] / (f)\) ，且 \(x\) 是 \(\mathbf{X}\) 在 \(\mathbf{A}\) 中的剩余类。对于每个 \(a \in \mathbf{K}\) 我们有 \(f(a) = \mathrm{N}_{\mathrm{A} / \mathrm{K}}(a - x)\) .
\item
  设K是一个域，A是次数为n的有限K-代数。对于每个与n互素的整数m，我们有 \(A^{*m} \cap K^{*} = K^{*m}\) ；换言之，如果 \(x \in K^{*}\) 并且如果存在 A 的一个元素 y 使得 \(x = y^{m}\) ，则存在 K 中的一个元素 z 使得 \(x = z^{m}\) .
\end{enumerate}

\BourbakiExerciseGroup

\begin{enumerate}
\def\labelenumi{\arabic{enumi})}
\tightlist
\item
  多项式 \(X^{2}-2\) 在 Q{[}X{]} 中不可约（参见 I，第 51 页，定理 7）；令 \(\alpha\) 是它在 Q 的一个代数闭包 \(\Omega\) 中的一个根。证明多项式 \(X^{2}-\alpha\) 在 \(E = Q(\alpha)\) 上不可约；设 \(\beta\) 是该多项式在 \(\Omega\) 中的一个根，并且令 \(F = E(\beta) = Q(\beta)\) 。证明 F 不是 Q 的拟伽罗瓦扩张；由 F 生成的 Q 的拟伽罗瓦扩张是什么？
\end{enumerate}

* 2) 设 E 为 Q 在 R 中的相对代数闭包（“实代数数域”）；若 u 是 E 的一个自同构，则 \(u(x) = x\) 对所有 \(x \in \mathbf{Q}\) 成立；证明若 \(y \in E\) 使得 \(y > 0\) ，则也有 \(u(y) > 0\) 。由此推出 \(u(y) = y\) 对所有 \(y \in E\) 成立. *

\begin{enumerate}
\def\labelenumi{\arabic{enumi})}
\setcounter{enumi}{2}
\item
  证明域 \(\mathbf{K}\) 的每个代数扩张，若它由一组在 \(\mathbf{K}\) 上次数为 2 的元素生成，则在 \(\mathbf{K}\) 上是拟伽罗瓦的.
\item
  设 K 为一个域，f 为 K{[}X{]} 中在 K 上可分且次数为 n 的不可约多项式，并设 \(\alpha_{i}\) ( \(1 \leqslant i \leqslant n\) ) 为其在 K 的一个代数闭包 \(\Omega\) 中的根。设 g 为 K{[}X{]} 中的一个多项式，且 h 为 K{[}X{]} 中多项式 \(f(g(\mathbf{X}))\) 的一个不可约因子。证明 h 的次数是 n 的倍数 rn，并且 h 与每个多项式 \(g(\mathbf{X}) - \alpha_{i}\) ( \(1 \leqslant i \leqslant n\) )（属于 \(\Omega[\mathbf{X}]\) ）恰好有 r 个公共根（考虑 h 的一个根在 K 上的共轭）。
\item
  \begin{enumerate}
  \def\labelenumii{\alph{enumii})}
  \tightlist
  \item
    设 K 为一个域， \(\Omega\) 为 K 的一个代数闭包，且 E、F 为 \(\Omega\) 的两个子扩张。证明 E 与 F 的每个复合扩张（V，第 12 页）都同构于形如 \((L_{w}, 1, w)\) 的一个复合扩张，其中 w 表示 F 到 \(\Omega\) 的一个子域上的 K-同构，1 是 E 的恒等映射，且 \(L_{w}\) 是 \(\Omega\) 中由 \(E \cup w(F)\) 生成的子域.
  \end{enumerate}
\end{enumerate}

\begin{enumerate}
\def\labelenumi{\alph{enumi})}
\setcounter{enumi}{1}
\item
  由 a) 推出，若 F 在 K 上的次数为有限的 n，则 E 与 F 的复合扩张中两两不同构的至多有 n 个。
\item
  设 E 是 K 的拟伽罗瓦扩张，F 是 E 的一个子扩张。证明 E 与 F 的每个复合扩张都同构于形如 \((E, 1, v)\) 的一个复合扩张，其中 v 是 F 到 E 的一个子域上的 K-同构。
\end{enumerate}

\begin{enumerate}
\def\labelenumi{\arabic{enumi})}
\setcounter{enumi}{5}
\tightlist
\item
  \begin{enumerate}
  \def\labelenumii{\alph{enumii})}
  \tightlist
  \item
    设 K 是特征指数为 p 的域，L 是 K 的一个扩张， \(\Omega\) 为 L 的一个代数闭包，以及 \(\bar{K}\) 是 K 在 \(\Omega\) 中的相对代数闭包。证明 \(\bar{K}\) 与 \(L^{p^{-\infty}}\) 在 \(E = \bar{K} \cap L^{p^{-\infty}}\) 上线性不相交。（设 \((a_{\lambda})\) 是向量空间 \(L^{p^{-\infty}}\) 在 E 上的一个基；考虑一个关系 \(\sum_{\lambda} c_{\lambda} a_{\lambda} = 0\) ，其中 \(c_{\lambda} \in \bar{K}\) 且非零的
  \end{enumerate}
\end{enumerate}

\(c_{\lambda}\) 的个数尽可能小，并把 \(\Omega\) 的一个 L-自同构作用于这个关系上。）

\begin{enumerate}
\def\labelenumi{\alph{enumi})}
\setcounter{enumi}{1}
\tightlist
\item
  由a)推出，若E和F是K的两个扩张，使得E在K上是代数的，且K在F中是相对代数闭的，则E和F的任意两个复合扩张都是同构的。
\end{enumerate}

\begin{enumerate}
\def\labelenumi{\arabic{enumi})}
\setcounter{enumi}{6}
\tightlist
\item
  设 K 是特征为 0 的域， \(\alpha\) , \(\beta\) , \(\xi\) 是 K 的一个代数闭包中的三个元素，使得 \([\mathbf{K}(\alpha):\mathbf{K}] = 2\) , \([\mathbf{K}(\beta):\mathbf{K}] = 3\) , \([\mathbf{K}(\xi):\mathbf{K}] = n > 1\) ；用 \(\alpha'\) （相应地 \(\beta'\) , \(\beta''\) ) 表示唯一的共轭元 \(\neq \alpha\) （ \(\alpha\) 的）（相应地，那些共轭元 \(\neq \beta\) （ \(\beta\) 的））.
\end{enumerate}

\begin{enumerate}
\def\labelenumi{\alph{enumi})}
\item
  证明： \([\mathbf{K}(\alpha, \xi): \mathbf{K}] = n\) 如果 \(\alpha \in \mathbf{K}(\xi)\) ，否则 \([\mathbf{K}(\alpha, \xi): \mathbf{K}] = 2n\) 。
\item
  证明： \([\mathbf{K}(\beta,\xi):\mathbf{K}]=n\) 如果 \(\beta\in\mathbf{K}(\xi)\) ；\([\mathbf{K}(\beta,\xi):\mathbf{K}]=2n\) 如果 \(\beta\in\mathbf{K}(\xi)\) 并且 \(\beta'\) , \(\beta''\) 中恰有一个属于 \(\mathbf{K}(\xi)\) ；最后 \([\mathbf{K}(\beta,\xi):\mathbf{K}]=3n\) 如果三个元素 \(\beta\) , \(\beta'\) , \(\beta''\) 中没有一个属于 \(\mathbf{K}(\xi)\) .
\item
  如果 \(d\) 是 \(\beta\) 的极小多项式的判别式，证明 \(\mathbf{K}(\beta, \beta') = \mathbf{K}(\beta, \sqrt{d})\) .
\end{enumerate}

\(d\) ) 证明 \(\mathbf{K}(\beta, \beta') = \mathbf{K}(\beta - \beta')\) .

\begin{enumerate}
\def\labelenumi{\alph{enumi})}
\setcounter{enumi}{4}
\item
  如果 \(n = 2\) 且 \(\alpha \notin \mathbf{K}(\xi)\) ，证明 \(\mathbf{K}(\alpha, \xi) = \mathbf{K}(\alpha + \xi)\) .
\item
  证明： \(\mathbf{K}(\alpha, \beta) = \mathbf{K}(\alpha + \beta)\) .
\end{enumerate}

\(g\) ) 证明 \(\mathbf{K}(\alpha ,\beta) = \mathbf{K}(\alpha \beta)\) ，除非 \(\alpha = aj\) 且 \(\beta^3 = b\) ，其中 \(a,b\) 是 \(\mathbf{K}\) 的元素且 \(j^3 = 1\) .

\begin{enumerate}
\def\labelenumi{\arabic{enumi})}
\setcounter{enumi}{7}
\tightlist
\item
  设 K 为无限域，E 为 K 的有限次代数扩张，且 \(\tau\) 是向量 K-空间 E 的一个自同构，使得对所有 \(x \in E\) , \(\tau(x)\) 在 K 上与 x 共轭。证明 \(\tau\) 是域 E 的一个 K-自同构。（设 N 为由 E 生成的 K 的拟伽罗瓦扩张，并设 \(\Gamma\) 为 N 的 K-自同构群， \(\Delta\) 为 \(\Gamma\) 的由所有 E-自同构组成的子群，并且 \(\sigma_{j}\) ( \(1 \leqslant j \leqslant r\) ) 是 \(\Delta\) 的左陪集的一个代表系。若 \((u_{i})_{1 \leqslant i \leqslant n}\) 是向量 K-空间 E 的一组基，考虑多项式
\end{enumerate}

\[
\mathrm{F} \left(\mathrm{X} _ {1}, \dots , \mathrm{X} _ {n}\right) = \prod_ {j = 1} ^ {r} \left(\sum_ {i = 1} ^ {n} \left(\tau \left(u _ {i}\right) - \sigma_ {j} \left(u _ {i}\right)\right) \mathrm{X} _ {i}\right).
\]

\BourbakiExerciseGroup

\begin{enumerate}
\def\labelenumi{\arabic{enumi})}
\tightlist
\item
  \begin{enumerate}
  \def\labelenumii{\alph{enumii})}
  \tightlist
  \item
    设 K 是一个域，f 是 K{[}X{]} 中的一个可分多项式，且 N 是 f 的一个分裂域。证明对于群 \(\operatorname{Gal}(N/K)\) 而言，要使它在 f 于 N 中的根集上可迁地作用，当且仅当 f 在 K{[}X{]} 中不可约.
  \end{enumerate}
\end{enumerate}

\begin{enumerate}
\def\labelenumi{\alph{enumi})}
\setcounter{enumi}{1}
\tightlist
\item
  进一步假设 \(f\) 在 \(\mathbf{K}[\mathbf{X}]\) 中不可约；设 \(x\) 是 \(f\) 在 \(\mathbf{N}\) 中的一个根。要使得不存在域 \(\mathbf{E}\) 满足 \(\mathbf{K} \subset \mathbf{E} \subset \mathbf{K}(x)\) 且异于 \(\mathbf{K}\) 和 \(\mathbf{K}(x)\) ，其充分必要条件是 \(\operatorname{Gal}(\mathbf{N} / \mathbf{K})\) 作为 \(f\) 在 \(\mathbf{N}\) 中的根集上的传递置换群应当是本原的（I，第 142 页，习题 13）。
\end{enumerate}

\begin{enumerate}
\def\labelenumi{\arabic{enumi})}
\setcounter{enumi}{1}
\tightlist
\item
  设 K 为一个域， \(f_{0}\) 为 K{[}X{]} 中一个次数为 n 的不可约且可分的多项式；设 \(\alpha_{i}\) ( \(1 \leqslant i \leqslant n\) ) 为它在 N 中的根，其中 N 是 \(f_{0}\) 的一个分裂域。设 F 为有理分式域 \(\mathrm{N}(\mathbf{X}_{1}, \mathbf{X}_{2}, \ldots, \mathbf{X}_{n})\) ，E 为 F 的子域 \(\mathrm{K}(\mathbf{X}_{1}, \mathbf{X}_{2}, \ldots, \mathbf{X}_{n})\) ；F 是 E 的伽罗瓦扩张，且 \(\operatorname{Gal}(\mathrm{E}/\mathrm{F})\) 典范同构于 \(\operatorname{Gal}(\mathrm{N}/\mathrm{K})\) (V，第 71 页，定理 5；V，第 154 页，习题 8)。证明若 \(\theta = \alpha_{1}\mathbf{X}_{1} + \alpha_{2}\mathbf{X}_{2} + \cdots + \alpha_{n}\mathbf{X}_{n}\) ，则 \(\mathrm{F} = \mathrm{E}(\theta)\) ，并且极小多项式 \(g_{1}\) （ \(\theta\) 在 E{[}T{]} 中的）是该多项式的一个不可约因子
\end{enumerate}

\[
f (T) = \prod_ {\sigma \in \mathfrak {S} _ {n}} \left(T - \alpha_ {1} X _ {\sigma (1)} - \alpha_ {2} X _ {\sigma (2)} - \dots - \alpha_ {n} X _ {\sigma (n)}\right).
\]

每个置换 \(\sigma \in \mathfrak{S}_n\) 定义了 E 的一个 K-自同构，仍记为 \(\sigma\) ，由条件 \(\sigma(X_i) = X_{\sigma(i)}\) （对 \(1 \leqslant i \leqslant n\) ）给定；证明 \(\operatorname{Gal}(N/K)\) 同构于 \(\mathfrak{S}_n\) 的由置换 \(\sigma\) （满足 \(\sigma(g_1) = g_1\) ）组成的子群。进一步地，若 \(f = g_1g_2 \ldots g_r\) 是 \(f\) 在 E{[}T{]} 中的不可约因子分解，证明对每个指标 \(j\) （满足 \(1 \leqslant j \leqslant r\) ）存在一个置换 \(\sigma_j \in \mathfrak{S}_n\) 使得 \(\sigma_j(g_1) = g_j\) .

\begin{enumerate}
\def\labelenumi{\arabic{enumi})}
\setcounter{enumi}{2}
\item
  设 N 是域 K 的无限次伽罗瓦扩张。证明伽罗瓦群 \(\operatorname{Gal}(N/K)\) 的基数大于 N 的子集所成集合的基数（利用 V，第 61 和 62 页）。由此推出存在 \(\operatorname{Gal}(N/K)\) 的非闭子群.
\item
  设 \(\mathbf{N}\) 是域 \(\mathbf{K}\) 的一个伽罗瓦扩张, \((\mathrm{E}_i)_{i\in I}\) 是 \(\mathbf{N}\) 的一族在 \(\mathbf{K}\) 上伽罗瓦的子扩张，使得： \(1^{\circ}\) 对每个指标 \(i\) ，如果 \(\mathbf{F}_i\) 是由 \(\mathbf{E}_j\) （ \(j\neq i\) ）生成的域，则 \(\mathbf{E}_i\cap \mathbf{F}_i = \mathbf{K}\) ; \(2^{\circ}\) \(\mathbf{N}\) 是由诸 \(\mathbf{E}_i\) 的并生成的.
\end{enumerate}

\begin{enumerate}
\def\labelenumi{\alph{enumi})}
\tightlist
\item
  证明： \(\mathbf{N}\) 同构于张量积 \(\otimes_{i\in I}\mathrm{E}_i\) （III，第470页）（定义一个
\end{enumerate}

这个张量积到N上的同构，利用V中第71页的定理5。

\begin{enumerate}
\def\labelenumi{\alph{enumi})}
\setcounter{enumi}{1}
\tightlist
\item
  证明： \(\operatorname{Gal}(\mathbf{N}/\mathbf{K})\) 同构于拓扑群的乘积 \(\operatorname{Gal}(\operatorname{E}_{i}/\operatorname{K})\) .
\end{enumerate}

\begin{enumerate}
\def\labelenumi{\arabic{enumi})}
\setcounter{enumi}{4}
\item
  证明每一个紧致完全不连通群都同构于某个域的伽罗瓦扩张的群Gal(N/K)（利用《一般拓扑学》III，第325页，习题3，化归到该群为有限群的乘积的情形，并同时使用习题4以及V，第59页，例5）。
\item
  设 N 是域 K 的一个伽罗瓦扩张，次数无限，且设 \(\rho, \sigma\) 是 \(\operatorname{Gal}(N/K)\) 的两个不同元素。如果 J 是由 \(x \in N\) （满足 \(\rho(x) \neq \sigma(x)\) ）组成的集合，证明 \(\mathrm{K}(\mathrm{J})\) 是 K 的无限次扩张。（注意，一个域不能是其两个子域的并集，除非它等于其中之一。）
\item
  设 \(\mathbf{N}\) 是域 \(\mathbf{K}\) 的一个伽罗瓦扩张，使得 \(\operatorname{Gal}(\mathbf{N} / \mathbf{K})\) 同构于对称群 \(\mathfrak{S}_n\) (V，第 59 页，例 5)，其中 \(n\) 为非素数的整数；把 \(\operatorname{Gal}(\mathbf{N} / \mathbf{K})\) 与 \(\mathfrak{S}_n\) 等同起来。设 \(\Gamma_1\) 为 \(\mathfrak{S}_n\) 的保持数字 1 不变的、阶为 \((n - 1)!\) 的子群，并设 \(\Gamma_2\) 为 \(\mathfrak{S}_n\) 的阶为 \(n\) 、由循环 \((1, 2, \ldots, n)\) 生成的循环子群（I，第 142 页，习题 12）。设 \(E_1, E_2\) 分别为子群 \(\Gamma_1, \Gamma_2\) 的不变量域。证明 \(E_1\) 与 \(E_2\) 在 \(\mathbf{K}\) 上线性不相交，不存在域 \(F\) （除 \(\mathbf{K}\) 和 \(E_1\) 之外）使得 \(K \subset F \subset E_1\) ，但存在域 \(F'\) （异于 \(E_2\) 和 \(N\) ）使得 \(E_2 \subset F' \subset N\) （利用 I，第 142 页习题 13 及 I，第 143 页习题 14）。
\item
  设 K 是一个域，N 是 K 的有限 n 次伽罗瓦扩张。
\end{enumerate}

\begin{enumerate}
\def\labelenumi{\alph{enumi})}
\item
  设 \(p\) 为一个整除 \(n\) 的素数，并令 \(p^r\) 为 \(p\) 的整除 \(n\) 的最高幂。证明存在子扩张 \(L_i\) （属于 \(N\) ， \(0 \leqslant i \leqslant r\) ）使得对于 \(0 \leqslant i \leqslant r - 1\) , \(L_i\) 包含 \(L_{i+1}\) , \(L_0 = N\) , \(L_i\) 是次数为 \(p\) 的 \(L_{i+1}\) 的伽罗瓦扩张，且 \([L_r: K]\) 不能被 \(p\) 整除（参见 I，第 77 页，定理 1 及 I，第 78 页，定理 2）。
\item
  设 \(x \in \mathbf{N}\) 使得 \(\mathbf{N}\) 由 \(x\) 及其共轭元生成。证明存在一个子扩张 \(\mathbf{E}\) （属于 \(\mathbf{N}\) ）使得 \(x \notin \mathbf{E}\) 且 \([\mathbf{N} : \mathbf{E}] = p'\) 。（考虑 Sylow \(p\) -群 \(\mathrm{H}_j\) （属于 \(\operatorname{Gal}(\mathbf{N} / \mathbf{K})\) ）以及由它们的并生成的（正规）子群 \(\mathbf{H}\) ；若 \(F_j\) （相应地 \(\mathbf{F}\) ）是 \(\mathrm{H}_j\) （相应地 \(\mathbf{H}\) ) 的不变子域，证明 \(x \notin \mathbf{F}\) 并由此推出存在一个指标 \(j\) 使得 \(x \notin F_j\) ；然后我们可以取 \(E = F_{j}\) 。) 得出结论：存在一个子扩张 \(L\) （属于 \(\mathbf{N}\) ）使得 \(L(x)\) 是次数为 \(p\) 的 \(L\) 的伽罗瓦扩张（利用 I，第 77 页，命题 12）。
\item
  反过来，设 \(\Omega\) 为 \(N\) 的一个代数闭包，并假设存在一个子扩张 \(L\) （属于 \(\Omega\) ）使得 \(L(x)\) 是次数为 \(p\) 的、关于 \(L\) 的伽罗瓦扩张。证明此时 \(p\) 整除 \(n\) 。（观察到 \(N \cap L(x) = (N \cap L)(x)\) 且 \((N \cap L)(x)\) 是次数为 \(p\) 的 \(N \cap L\) 的伽罗瓦扩张.)
\end{enumerate}

\begin{enumerate}
\def\labelenumi{\arabic{enumi})}
\setcounter{enumi}{8}
\tightlist
\item
  设 K 是一个特征为 \(p\) 的域, \(\Omega\) 为 K 的一个代数闭扩张， \(x\) 为 \(\Omega\) 的一个元素，使得 \(x \notin K\) ，且 M 为 \(\Omega\) 的不包含 \(x\) 的子扩张所成集合的一个极大元（V，第 148 页，习题 9），因此 \(\Omega\) 是 M 的代数扩张（同上引）。a) 证明：若 M 不是完全域， \(\Omega\) 是 M 的一个 \(p\) -根式扩张，即诸域 \(M(x^{p^{-n}})\) （ \(n \geqslant 0\) ）的并（注意对于每个元素 \(y \in \Omega\) （在 M 上为 \(p\) -根式的），域 \(M(y)\) 包含 \(M(x)\) ）；由此推出我们必然有 \([M(x):M] = p\) ，因为 \(\Omega\) 的每个在 M 上可分的元素必然属于 M）。
\end{enumerate}

\begin{enumerate}
\def\labelenumi{\alph{enumi})}
\setcounter{enumi}{1}
\tightlist
\item
  证明若 M 是完全域， \(\mathbf{M}(x)\) 是 M 的伽罗瓦扩张，次数为素数 q，并且对于每个子扩张 N （ \(\Omega\) 的、在 M 上为有限次伽罗瓦的），其次数 \([N:M]\) 是 q 的幂。（取 q 为整除 \([\mathbf{M}(x):\mathbf{M}]\) 的一个素数；注意 M 的每个不同于 M 且包含于 \(\Omega\) 的扩张都包含 \(\mathbf{M}(x)\) ，并利用习题 8 a)。）
\end{enumerate}

\begin{enumerate}
\def\labelenumi{\arabic{enumi})}
\setcounter{enumi}{9}
\tightlist
\item
  设 \(\mathbf{K}\) 是一个完全域，使得对于每个整数 \(n > 1\) 存在一个伽罗瓦扩张 \(\mathbf{N}\) （ \(\mathbf{K}\) 的）使得 \(\operatorname{Gal}(\mathbf{N} / \mathbf{K})\) 同构于交错群 \(\mathfrak{A}_n\) ；例如，对于每个特征为 0 的域 \(k\) ，域 \(\mathbf{K} = k(\mathbf{X}_n)_{n \in \mathbb{N}}\) （无穷多个未定元上的有理分式域）具有此性质（V，第 59 页，例 5）。
\end{enumerate}

\begin{enumerate}
\def\labelenumi{\alph{enumi})}
\item
  设 A 为一个由 \textgreater{} 1 的整数组成的有限集，且设 \(\Omega\) 为 K 的一个代数闭包。一个元素 \(x \in \Omega\) 称为 A-可构造的，如果存在 K 的扩张组成的一个递增序列 \(K = E_{0} \subset E_{1} \subset \cdots \subset E_{r} \subset \Omega\) ，使得每个次数 \([E_{j}: E_{j-1}]\) （对于 \(1 \leqslant j \leqslant r\) ）都属于 A，且 \(x \in E_{r}\) 。证明存在一个极大元 L ，于子扩张 \(E \subset \Omega\) （其中每个 \(x \in E\) 是 A-可构造的）所成的集合中。对于每个 \(x \notin L\) 证明次数 \([L(x): L]\) 属于 N - A。
\item
  设 \(a\) 是 \(\mathbf{A}\) 的最大元，并且令 \(n\) 为一个满足 \(n \geqslant \max (a + 1,5)\) 的整数。证明存在一个扩张 \(\mathbf{M}\) （ \(\mathbf{L}\) 的）使得 \([\mathbf{M}:\mathbf{L}] = n\) 。（首先观察到，存在一个不可约多项式 \(f \in \mathbf{K}[\mathbf{K}]\) ，其次数为 \(n\) ，使得如果 \(\mathbf{N}\) 是 \(f\) 的分裂域， \(\operatorname{Gal}(\mathbf{N} / \mathbf{K})\) 同构于 \(\mathfrak{A}_n\) ；设 \(y \in \Omega\) 为 \(f\) 在 \(\Omega\) 中的一个根并且令 \(E = K(y)\) 。证明我们必然有 \(E \cap L = K\) ，因此 \(L(E)\) 是所要求的域。用反证法论证，注意到 \(E \cap L\) 的元素是 A-可构造的，并且关系 \(E \cap L \neq K\) 将蕴含一个正规子群 \(\Gamma\) （ \(\mathfrak{A}_n\) 的，满足 \(1 < (\mathfrak{A}_n:\Gamma) < n\) ）的存在性，这将与 \(\mathfrak{A}_n\) 的单性相矛盾（I，第 143 页，习题 16）。
\end{enumerate}

\begin{enumerate}
\def\labelenumi{\arabic{enumi})}
\setcounter{enumi}{10}
\tightlist
\item
  \begin{enumerate}
  \def\labelenumii{\alph{enumii})}
  \tightlist
  \item
    设 E 是域 K 的有限次可分扩张，N 是由 E 生成的 K 的伽罗瓦扩张。若 \(\alpha\) 是 N 的一个元素，其共轭元构成 N 在 K 上的正规基，且 \(\beta = \operatorname{Tr}_{\mathbb{N}/\mathbb{E}}(\alpha)\) ，证明 \(\operatorname{E} = \operatorname{K}(\beta)\) .
  \end{enumerate}
\end{enumerate}

\begin{enumerate}
\def\labelenumi{\alph{enumi})}
\setcounter{enumi}{1}
\tightlist
\item
  设 N 是 K 的有限次伽罗瓦扩张，且令 \(\alpha\) 为 N 的一个元素，其共轭元在 K 上构成一个正规基。若 L 是 N 的一个在 K 上伽罗瓦的子扩张，且 \(\beta = \mathrm{Tr}_{\mathbb{N}/\mathbb{L}}(\alpha)\) ，证明 \(\beta\) 的共轭元构成 L 在 K 上的一个正规基。
\item
  设 L、M 是 K 的两个有限次伽罗瓦扩张，使得 \(L \cap M = K\) ；设 \(\alpha\) （相应地 \(\beta\) ）是 L（相应地，M）中的一个元素，其在 K 上的共轭元构成 L（相应地，M）在 K 上的一个正规基；证明 \(\alpha\beta\) 的在 K 上的共轭元构成 \(\mathrm{K}(\mathrm{L} \cup \mathrm{M})\) 的一个正规基.
\end{enumerate}

\begin{enumerate}
\def\labelenumi{\arabic{enumi})}
\setcounter{enumi}{11}
\tightlist
\item
  设 K 为无限域，N 为 K 的 n 次伽罗瓦扩张， \(f \in K[X]\) 为 N 中一个元素 \(\theta\) 的极小多项式，使得 \(\mathbf{N} = \mathbf{K}(\theta)\) 。群 \(\Gamma = \operatorname{Gal}(N/K)\) 可以作用在 N{[}X{]} 中的多项式上，方式是作用于这些多项式的系数。
\end{enumerate}

\begin{enumerate}
\def\labelenumi{\alph{enumi})}
\tightlist
\item
  设 \(g \in \mathbb{N}[\mathbf{X}]\) 为次数 \(\leqslant n - 1\) 的、使得 \(g(\theta) = 1\) , \(g(\sigma(\theta)) = 0\) （对所有 \(\sigma \neq 1\) ， \(\Gamma\) 中的）的多项式。对于每个子群 \(\Delta\) （ \(\Gamma\) 的），令 \(p_{\Delta} \in \mathbb{N}[\mathbf{X}]\) 为多项式 \(\det(\sigma\tau(g))_{\sigma \in \Delta, \tau \in \Delta}\) 。证明
\end{enumerate}

\[
p _ {\Delta} ^ {2} \equiv \sum_ {\sigma \in \Delta} \sigma (g) \mod f
\]

并推导出 \(p_{\Delta} \neq 0\) .

\begin{enumerate}
\def\labelenumi{\alph{enumi})}
\setcounter{enumi}{1}
\tightlist
\item
  推导出存在一个元素 \(\alpha \in K\) ，使得对于每个子群 \(\Delta\) （ \(\Gamma\) 的）， \(g(\alpha)\) 在子域 \(k(\Delta)\) 上的共轭元构成 N 在 \(k(\Delta)\) 上的一个正规基.
\end{enumerate}

\begin{enumerate}
\def\labelenumi{\arabic{enumi})}
\setcounter{enumi}{12}
\item
  设 E 是域 K 的一个代数扩张， \(E_{s}\) 为 K 的包含在 E 中的最大可分扩张， \(E_{r}\) 为 K 的包含在 E 中的最大 p-根式扩张。在 K 的一个代数闭包中，设 N 是由 E 生成的 K 的拟伽罗瓦扩张；则包含在 N 中的 K 的最大可分扩张是拟伽罗瓦扩张 \(N_{s}\) （由 \(E_{s}\) 生成的）。要使 \(E_{r}\) 是包含在 N 中的 K 的最大 p-根式扩张，必要且充分条件是 E 在 \(E_{r}\) 上可分（参见 V，第 152 页，习题 3）。
\item
  设 \(\mathbf{E}\) 是域 \(\mathbf{K}\) 的一个代数扩张, \(\Gamma\) 为 \(\mathbf{K}\) -自同构组成的群（ \(\mathbf{E}\) 的），且 \(\mathbf{S} \subset \mathbf{E}\) 为 \(\Gamma\) 的不变量域.
\end{enumerate}

\begin{enumerate}
\def\labelenumi{\alph{enumi})}
\item
  证明 E 在 K 上拟伽罗瓦的充分必要条件是 S 是 K 的 p-根式扩张。
\item
  设 \(S_{s}\) 为包含在 S 中的 K 的最大可分扩张。证明 \(S_{s}\) 是满足以下条件的域 F 中的最小者： \(K \subset F \subset E\) 且 E 是 F 的拟伽罗瓦扩张。
\item
  设 \(E_{s}\) 为包含在 E 中的 K 的最大可分扩张；证明 \(\mathrm{E} = \mathrm{S}(\mathrm{E}_{s})\) （注意，E 的 K-自同构中除恒等映射外，没有任何一个会保持 \(\mathrm{S}(\mathrm{E}_{s})\) 的所有元素不变）。
\end{enumerate}

\begin{enumerate}
\def\labelenumi{\arabic{enumi})}
\setcounter{enumi}{14}
\item
  设 E 是域 K 的有限次 n 的伽罗瓦扩张，其伽罗瓦群 G 是幂零的。设 P 为整除 n 的素数集合。则存在 E 的伽罗瓦子扩张 \(E_{p}\) （ \(p \in P\) ）使得 \(\operatorname{Gal}(E_{p}/K)\) 对每个 \(p \in P\) 都是 p-群，并且使得典范同态 \(\otimes E_{p} \to E\) 是一个同构。
\item
  设 \(a\) 是一个整数，并令 \(x_1, x_2, x_3, x_4\) 是 \(P(X) = X^4 - aX - 1\) 在 \(\mathbf{Q}\) 的一个代数闭包中的根。我们把伽罗瓦群 \(G\) （ \(P\) 的）与 \(\{x_1, \ldots, x_4\}\) 上的一个置换群等同起来。假设 \(Q(x_1)\) 不包含二次子扩张（参见 \(V\) ，第 154 页，习题 1）。
\end{enumerate}

\begin{enumerate}
\def\labelenumi{\alph{enumi})}
\item
  证明： \(G = S_{4}\) 或 \(A_{4}\) （利用 V 章第 154 页习题 1：P 不可约，因此 G 传递作用；进一步 \(\mathbf{Q}(x_{1})\) 不包含二次子扩张，因此 \(x_{1}\) 的稳定化子 H 不包含在 G 的指数为 2 的子群中，而 \([G:H] = 4\) ；由此推出 G 不是 2-群，因此其阶可被 3 整除。）
\item
  设 \(\delta = \prod_{1\leqslant i < j\leqslant 4}(x_i - x_j)\) ，且 \(d = \delta^2\) 。证明 \(d = 4^4 -3^3 a^4\) ；进一步有 \(s(\delta) = \varepsilon_s\delta\) ，对
\end{enumerate}

所有 \(s \in G\) 成立。由此推出 \(G = \mathfrak{A}_4\) 如果 \(\delta \in Q\) ，而 \(G = \mathfrak{S}_4\) 否则。证明 \(\delta\) 不是有理的。（因此我们构造了 \(Q\) 上的、具有伽罗瓦群 \(\mathfrak{S}_4\) 的 4 次方程的一个无穷族。）

* 17) 设 \(f\) 是 \(\mathbf{Q}[\mathbf{X}]\) 中的素数次数 \(p\) 的不可约多项式，在 \(\mathbf{C} - \mathbf{R}\) 中恰好有两个根。则伽罗瓦群 \(G\) （ \(f\) 的）同构于 \(\mathfrak{S}_p\) 。（G 包含一个阶为 \(p\) 的元素和一个对换。）*

\begin{enumerate}
\def\labelenumi{\arabic{enumi})}
\setcounter{enumi}{17}
\tightlist
\item
  设 K 为一个域，F 为有理分式域 \(\mathrm{K}(\mathrm{T})\) 。群 \(\mathrm{GL}_{2}(\mathrm{K})\) 按法则 \((\gamma f)(\mathrm{T}) = f\left(\frac{a\mathrm{T} + b}{c\mathrm{T} + d}\right)\) 作用于 F，其中 \(\gamma = \begin{pmatrix} a & b \\ c & d \end{pmatrix} \in \mathrm{GL}_{2}(\mathrm{K})\) 且 \(f \in F\) 。中心 \(K^{*}\) .I （即 \(\mathrm{GL}_{2}(\mathrm{K})\) 的单位元）平凡地作用，因此通过取商，我们得到 \(\mathrm{PGL}_{2}(\mathrm{K}) = \mathrm{GL}_{2}(\mathrm{K}) / \mathrm{K}^{*}\) .I 到 K 的扩张 F 的自同构群内的一个同态。这个同态是双射的（V，第 148 页，习题 11）。
\end{enumerate}

\begin{enumerate}
\def\labelenumi{\alph{enumi})}
\tightlist
\item
  设 H 是 \(\mathrm{PGL}_{2}(\mathbf{K})\) 的一个 h 阶有限子群，并设 E 为 H 的不变量域，从而 \(F = E(T)\) 是 E 的伽罗瓦扩张，其伽罗瓦群为 H（V，第 66 页，定理 3）。设
\end{enumerate}

\[
\mathrm{H} (\mathrm{X}) = \prod_ {s \in \mathrm{H}} (\mathrm{X} - s (\mathrm{T})) = \sum_ {i = 0} ^ {h} (- 1) ^ {i} \mathrm{S} _ {i} (\mathrm{T}) \mathrm{X} ^ {h - i}
\]

为 T 在 E 上的极小多项式。对于每个满足 \(0 \leqslant i \leqslant h\) 的 i，或者 \(S_{i}(T) \in K\) ，或者 \(E = K(S_{i}(T))\) 。（利用 V，第 157 页，习题 11 a)。）

\begin{enumerate}
\def\labelenumi{\alph{enumi})}
\setcounter{enumi}{1}
\item
  对于每个 \(s \in H\) ，令 \(\begin{pmatrix} a_{s} & b_{s} \\ c_{s} & d_{s} \end{pmatrix}\) 为它在 \(GL_{2}(K)\) 中的一个代表元，并且令 \(H_{0}\) 为 H 中由满足 \(c_{s} = 0\) 的那些 s 组成的子群；设 \(h_{0} = CardH_{0}\) 并令 \(R^{H}(T) = S_{h_{0}}(T)\) 。证明 \(E = K(R^{H}(T))\) 且 \(R^{H}(T) = P(T)/Q(T)^{h_{0}}\) ，其中 P 是次数为 h 的多项式，且 \(Q(T) = \prod_{s} (c_{s}T + d_{s})\) ，其中 s 遍历陪集 \(H_{0}s\) （异于 \(H_{0}\) 的）的一个代表系；因此 Q 的次数为 \((h/h_{0}) - 1\) .
\item
  假设 \(K^{*}\) 包含一个有限阶 n 的子群 \(\mu_{n}\) 。设 \(C_{n}\) 为（在 \(PGL_{2}(K)\) 中的）\(GL_{2}(K)\) 的由矩阵 \(\begin{pmatrix} a & 0 \\ 0 & 1 \end{pmatrix}, a \in \mu_{n}\) 组成的子群的像。则 \(C_{n}(X) = X^{n} - T^{n}\) 且 \(R^{C_{n}}(T) = (-1)^{n+1}T^{n}\) .
\item
  设 \(w\) 为 \(\mathrm{PGL}_2(\mathbf{K})\) 中由矩阵 \(\begin{pmatrix} 0 & 1 \\ 1 & 0 \end{pmatrix}\) 表示的元素。则 \(w^2 = 1\) ，且 \(wsw^{-1} = s^{-1}\) （若 \(s\) 由一个对角矩阵表示）。子群 \(\mathrm{D}_n\) （ \(\mathrm{PGL}_2(\mathbf{K})\) 的、由 \(\mathrm{C}_n\) 与 \(w\) 生成的）因此是 \(2n\) 阶的。我们有 \(\mathrm{D}_n(\mathbf{X}) = (\mathbf{X}^n - \mathrm{T}^n)\left(\mathbf{X}^n - \left(\frac{1}{\mathrm{T}}\right)^n\right) = \mathbf{X}^{2n} + (-1)\mathbf{R}^{\mathrm{D}_n(\mathrm{T})} + 1\) ，其中 \(\mathbf{R}^{\mathrm{D}_n(\mathrm{T})} = -((-T)^n + (-T)^{-n})\) .
\item
  设 S 为 \(\mathrm{PGL}_2(\mathbf{K})\) 中由矩阵 \(\left( \begin{array}{cc}0 & -1\\ 1 & 1 \end{array} \right)\) 的像生成的 3 阶群。我们有 \(S(X) = (X - T)\left(X - \frac{1}{1 - T}\right)\left(X - \frac{T - 1}{T}\right)\) 且 \(R^s (T) = \frac{T^3 - 3T + 1}{T^2 - 1}\) .
\item
  设 K 是一个含有 q 个元素的有限域。令 U 为（在 \(\mathrm{PGL}_{2}(\mathbf{K})\) 中的）由所有矩阵 \(\begin{pmatrix}1 & 0 \\ b & 1\end{pmatrix}\) , \(b \in K\) 构成的群的像。则
\end{enumerate}

\[
\mathrm{U} (\mathrm{X}) = (\mathrm{X} - \mathrm{T}) ^ {q} - (\mathrm{X} - \mathrm{T}) = \mathrm{X} ^ {q} - \mathrm{X} - (\mathrm{T} ^ {q} - \mathrm{T}) \quad \text { and } \quad \mathrm{R} ^ {\mathrm{U}} (\mathrm{T}) = \mathrm{T} ^ {q} - \mathrm{T}.
\]

设 B 是由 \(C_{q-1}\) 与 U 生成的群；其阶为 \(q(q-1)\) 。我们有

\[
\mathrm{B} (\mathrm{X}) = \left(\mathrm{X} ^ {q} - \mathrm{X}\right) ^ {q - 1} - \left(\mathrm{T} ^ {q} - \mathrm{T}\right) ^ {q - 1} \quad \text { and } \quad \mathrm{R} ^ {\mathrm{H}} (\mathrm{T}) = \left(\mathrm{T} ^ {q} - \mathrm{T}\right) ^ {q - 1}.
\]

\begin{enumerate}
\def\labelenumi{\arabic{enumi})}
\setcounter{enumi}{18}
\item
  设 E 是域 K 的伽罗瓦扩张，伽罗瓦群为 G。对每个 \(x \in E\) 用 Gx 表示 x 在 E 中的共轭元素的集合；这是一个有限的 G-集。设 x，y 为 E 中的两个元素。要使 G-集 Gx 和 Gy 同构（即存在一个双射 \(Gx \rightarrow Gy\) ，与 G 在这两个集合上的作用相容），其必要且充分条件是扩张 \(\mathrm{K}(x)\) 与 \(\mathrm{K}(y)\) 共轭。
\item
  设 E 是域 K 的一个代数扩张，使得 K{[}X{]} 中每个非常数多项式在 E 中至少有一个根。证明 E 是代数闭的。（在 E 的代数闭包中——从而也是在 K 的代数闭包中——设 F 是 K 的有限次扩张。要证明 \(F \subset E\) ，可归结为 F 是拟伽罗瓦扩张的情形，然后利用 V 第 76 页命题 13，再归结为 F 是伽罗瓦扩张或 p-根式扩张的情形。）
\item
  设 \(f(\mathbf{X}) = \mathbf{X}^{n} + a_{1}\mathbf{X}^{n-1} + \cdots + a_{n}\) 是 Z{[}X{]} 中的多项式，且 p 为素数。假设 \(a_{i} \equiv 0 \pmod{p}\) , \(1 \leqslant i \leqslant n\) ，且 \(a_{n} \not\equiv 0 \pmod{p^{2}}\) 。证明 \(f(\mathbf{X})\) 在 Q{[}X{]} 中不可约。（设 \(f(\mathbf{X}) = \mathrm{P}(\mathbf{X}) \mathrm{Q}(\mathbf{X})\) 是 Q{[}X{]} 中的一个分解。证明：若 P 和 Q 为首一多项式，则它们具有整数系数。（模 p 约化。）
\end{enumerate}

* 22) 设 \(m\) 是一个偶数 \(>0\) , \(r\) 为一个整数 \(\geqslant 3\) ，且 \(n_1 < \cdots < n_{r-2}\) 为 \(r-2\) 个偶数。令 \(g(\mathbf{X}) = (\mathbf{X}^2 + m)(\mathbf{X} - n_1) \ldots (\mathbf{X} - n_{r-2})\) .

\begin{enumerate}
\def\labelenumi{\alph{enumi})}
\item
  令 \(f(\mathbf{X}) = g(\mathbf{X}) - 2\) 。证明 \(f(\mathbf{X})\) 至少有 \(r - 2\) 个实根。计算 \(f\) 在 \(\mathbf{C}\) 中的根的平方和。由此推出，对于 \(m \geqslant \sum n_i^2 / 2\) , \(f(\mathbf{X})\) 是一个具有整数系数的多项式，它有 \(r - 2\) 个实根和两个非实的共轭复根。
\item
  证明： \(f(\mathbf{X})\) 在 \(\mathbf{Q}[\mathbf{X}]\) 中不可约。（写 \(f(\mathbf{X}) = \mathbf{X}' + \sum_{k=1}^{r} a_k \mathbf{X}'^{-k}\) ；证明
\end{enumerate}

\(a_{k}\) 是偶数且 \(a_{r}\) 不能被4整除；应用习题21。）

\begin{enumerate}
\def\labelenumi{\alph{enumi})}
\setcounter{enumi}{2}
\item
  证明当 \(r\) 为素数且 \(m \geqslant \sum n_i^2 / 2\) 时，分裂域的伽罗瓦群 \(G\) （ \(f\) 的）同构于 \(\mathfrak{S}_r\) （参见 V，第 158 页，习题 17）。
\item
  证明当 \(r = 4\) 且 \(m \geqslant \sum n_i^2 / 2\) 时，分裂域在 \(\mathbf{Q}\) 上的次数等于 8 或 24，并且它等于 8 当且仅当 \(n_1 = -n_2\) ，或者也当且仅当 \(f(\mathbf{X})\) 具有形式 \(\mathrm{P}(\mathbf{X}^2)\) ，其中 \(\mathbf{P}\) 是一个次数为 2 的多项式。*
\end{enumerate}

\begin{enumerate}
\def\labelenumi{\arabic{enumi})}
\setcounter{enumi}{22}
\tightlist
\item
  设 \(k\) 是一个特征为2的域。
\end{enumerate}

\begin{enumerate}
\def\labelenumi{\alph{enumi})}
\tightlist
\item
  设 \(A_{1}, \ldots, A_{n}\) 为未定元，K 为域 \(k(\mathbf{A}_{1}, \ldots, \mathbf{A}_{n})\) ，E 为多项式 \(P = X^{n} + A_{1}X^{n-1} + \cdots + A_{n} \in K[X]\) 的一个分裂域，且 \(X_{1}, \ldots, X_{n}\) 为 P 在 E 中的根，因此 \(E = k(X_{1}, \ldots, X_{n})\) （V，第 21 页）。我们将 E/K 的伽罗瓦群与对称群 \(S_{n}\) （通过变量 \(X_{1}, \ldots, X_{n}\) 的置换而作用）等同起来（V，第 59 页）。设
\end{enumerate}

\[
\beta \left(\mathrm{X} _ {1}, \dots , \mathrm{X} _ {n}\right) = \sum_ {i <   j} \mathrm{X} _ {i} / \left(\mathrm{X} _ {i} + \mathrm{X} _ {j}\right) \in \mathrm{E} \text { and } \mathrm{C} \left(\mathrm{A} _ {1}, \dots , \mathrm{A} _ {n}\right) = \sum_ {i <   j} \left(\mathrm{X} _ {i} \mathrm{X} _ {j}\right) / \left(\mathrm{X} _ {i} ^ {2} + \mathrm{X} _ {j} ^ {2}\right) \in \mathrm{K}.
\]

则 \(\beta \notin K\) 且

\[
\boldsymbol {\beta} ^ {2} + \boldsymbol {\beta} + \mathbf {C} = 0.
\]

我们有 \(\sigma(\beta) = \beta\) （对 \(\sigma \in \mathfrak{A}_n\) ），且 \(\sigma(\beta) = \beta + 1\) （对 \(\sigma \in \mathfrak{S}_n - \mathfrak{A}_n\) ）。由此推出 \(K(\beta)\) 是 \(E/K\) 的唯一的二次子扩张.

\begin{enumerate}
\def\labelenumi{\alph{enumi})}
\setcounter{enumi}{1}
\tightlist
\item
  设 \(D \in Z[T_{1}, \ldots, T_{n}]\) 为多项式 \(X^{n} + T_{1}X^{n-1} + \cdots + T_{n}\) 的判别式。证明存在 Q, R， \(S \in Z[T_{1}, \ldots, T_{n}]\) ，满足
\end{enumerate}

\[
\mathbf {D} = \mathbf {Q} ^ {2} - 4 \mathbf {R} + 8 \mathbf {S}
\]

且

\[
\mathrm{C} = \mathrm{C} (\mathrm{A} _ {1}, \dots , \mathrm{A} _ {n}) = \frac {\mathrm{R} (\mathrm{A} _ {1} , \dots , \mathrm{A} _ {n})}{\mathrm{Q} ^ {2} (\mathrm{A} _ {1} , \dots , \mathrm{A} _ {n})} = \frac {\mathrm{R} (\mathrm{A} _ {1} , \dots , \mathrm{A} _ {n})}{\mathrm{D} (\mathrm{A} _ {1} , \dots , \mathrm{A} _ {n})}.
\]

对于 \(\mathbf{Z}[T_1, \ldots, T_n]\) 的元素 U, V, W 的每一个满足 \(D = U^2 - 4V + 8W\) 的三元组，存在 \(M \in \mathbf{Z}[A_1, \ldots, A_n]\) 使得

\[
C = \frac {V (A _ {1} , \dots , A _ {n})}{D (A _ {1} , \dots , A _ {n})} + \left(\frac {M}{D (A _ {1} , \dots , A _ {n})}\right) ^ {2} + \frac {M}{D (A _ {1} , \dots , A _ {n})}.
\]

\begin{enumerate}
\def\labelenumi{\alph{enumi})}
\setcounter{enumi}{2}
\tightlist
\item
  设 \(f = X^{n} + a_{1}X^{n-1} + \cdots + a_{n}\) 是 k{[}X{]} 的一个可分多项式，且 G 是一个分裂扩张的伽罗瓦群。证明 \((a_{1}, \ldots, a_{n})\) 在 C 中是可代入的，并且以下条件等价
\end{enumerate}

\begin{enumerate}
\def\labelenumi{(\roman{enumi})}
\item
  \(G\) 的每个元素诱导 \(f\) 的根的一个偶置换;
\item
  存在 \(\alpha \in k\) 使得 \(C(a_{1},\ldots ,a_{n}) = \alpha^{2} + \alpha\) .
\end{enumerate}

* \(d\) ) 今后假设 \(k\) 是有限的，并任意选取 \(\mathbf{Z}[T_1, \ldots, T_n]\) 的元素 U、V、W 以满足 \(b\) ) 中所述的条件.

证明 \(c\) ) 的 (i) 和 (ii) 与以下条件等价：

\begin{enumerate}
\def\labelenumi{(\roman{enumi})}
\setcounter{enumi}{2}
\tightlist
\item
  \[
\operatorname{Tr} _ {k / \mathbf {F} _ {2}} (\mathbf {C}) = 0;
\]
\item
  \[
\mathrm{Tr} _ {k / \mathbf {F} _ {2}} \left(\frac {\mathbf {V} (a _ {1} , \dots , a _ {n})}{\mathbf {D} (a _ {1} , \dots , a _ {n})}\right) = 0.
\]
\end{enumerate}

\begin{enumerate}
\def\labelenumi{(\alph{enumi})}
\setcounter{enumi}{21}
\tightlist
\item
  \(f\) 的不可约因子的个数与 \(n\) 模 2 同余（使用 V，第 165 页，习题 22）。*
\end{enumerate}

\BourbakiExerciseGroup

\begin{enumerate}
\def\labelenumi{\arabic{enumi})}
\item
  设 \(\mathbf{K}\) 是特征为 \(p\) 的域， \(n\) 为一个不是 \(p\) 的倍数的整数， \(\mathbf{R}_n(\mathbf{K})\) 为 \(n\) 次单位根域，以及 \(\mathbf{E}\) 为 \(\mathbf{R}_n(\mathbf{K})\) 的一个循环扩张，其次数为 \(n\) （在 \(\mathbf{K}\) 上）；\(\mathbf{E}\) 因此是多项式 \(X^{n}-\alpha\in\mathbb{R}_{n}(\mathbf{K})[\mathbf{X}]\) 的根域（V，第 89 页，定理 4）。要使 E 成为 K 的伽罗瓦扩张，必要且充分的条件是：对每个自同构 \(\tau\in\operatorname{Gal}\left(\mathbb{R}_{n}(\mathbf{K})/\mathbf{K}\right)\) ，存在一个整数 r\textgreater0 和一个元素 \(\gamma\in\mathbb{R}_{n}(\mathbf{K})\) 使得 \(\tau(\alpha)=\gamma^{n}\alpha^{r}\) .
\item
  \begin{enumerate}
  \def\labelenumii{\alph{enumii})}
  \tightlist
  \item
    设 E 是多项式 \(X^{3}-2\in Q[X]\) 的一个分裂扩张。域 E 是 \(\mathrm{R}_{3}(\mathbf{Q})\) 的 3 次循环扩张，但不包含 Q 的 3 次循环扩张。b) 域 \(\mathrm{E}=\mathrm{R}_{9}(\mathbf{Q})\) 是 \(\mathrm{R}_{3}(\mathbf{Q})\) 的 3 次循环扩张，并且包含 Q 的 3 次循环扩张。
  \end{enumerate}
\end{enumerate}

\begin{enumerate}
\def\labelenumi{\alph{enumi})}
\setcounter{enumi}{2}
\tightlist
\item
  在 V，第 154 页，习题 7 的假设和记号下，证明 \(\mathrm{K}(\alpha,\beta)\) 在 K 上为伽罗瓦扩张的充要条件是 d 是 \(\mathrm{K}(\alpha)\) 中的平方；\(\mathrm{K}(\beta)\) 在 K 上为循环扩张的充要条件是 d 是 K 中的平方。
\end{enumerate}

* 3) 证明：对每个整数 \(n > 1\) 都存在一个次数为 \(n\) 的、关于域 \(\mathbf{Q}\) 的循环扩张。（化归到 \(n\) 是一个素数 \(q\) 的情形，并利用乘法群 \((\mathbf{Z} / q^m\mathbf{Z})^*\) 的结构（VII，第 13 页）。*

\begin{enumerate}
\def\labelenumi{\arabic{enumi})}
\setcounter{enumi}{3}
\tightlist
\item
  设 \(\mathbf{K}\) 是特征为 \(p, q\) 一个素数 \(\neq p\) ，使得 \(\mathbf{K}\) 包含 \(q\) 次单位根，并设 \(\zeta\) 是一个本原 \(q\) 次单位根。
\end{enumerate}

\begin{enumerate}
\def\labelenumi{\alph{enumi})}
\item
  设 E 是 K 的循环扩张，次数为 \(q^e\) ，且 \(\sigma\) 是 E 的一个生成 Gal(E/K) 的 K-自同构。设 F 是 K 上次数为 \(q^{e-1}\) 的域，使得 \(K \subset F \subset E\) ；存在 \(\theta \in E\) ，为某个不可约多项式 \(X^q - \alpha\) （在 F{[}X{]} 中）的一个根，使得 \(E = F(\theta)\) 且 \(\theta^{\sigma^m} = \zeta\theta\) （其中 \(m = q^{e-1}\) ）。证明同样有 \(E = K(\theta)\) 且 \(\theta^\sigma = \beta\theta\) ，其中 \(\beta \in F\) 使得 \(\alpha^{\sigma-1} = \beta^q\) 且 \(N_{F/K}(\beta) = \zeta\) （观察到 \(\theta^\sigma = \beta\theta^k\) ，其中 \(0 < k < q\) 且 \(\beta \in F\) ，由习题 1；通过计算 \(\theta^{\sigma^{mq}}\) ，证明 \(k^{mq} - 1 \equiv 0 \pmod{q}\) ）并推断出 \(k = 1\) ).
\item
  反之，设 F 是 K 的一个次数为 \(q^{e-1}\) 的循环扩张，其中 \(e \geqslant 2\) ，并且令 \(\sigma\) 是 F 的一个生成 Gal(F/K) 的 K-自同构。若存在一个元素 \(\beta \in F\) 使得 \(N_{F/K}(\beta) = \zeta\) ，并且 \(\alpha \in F\) 使得 \(\alpha^{\sigma-1} = \beta^q\) ，证明对每一个 \(c \in K^*\) ，多项式 \(X^q - c\alpha\) 在 F{[}X{]} 中不可约。若 \(\theta\) 是该多项式在 F 的一个代数闭包 \(\Omega\) 中的一个根，证明存在一个 K-同态 \(\bar{\sigma}\) ，它把 \(E = F(\theta)\) 映到 \(\Omega\) 内，扩张 \(\sigma\) ，并且满足 \(\theta^{\sigma} = \beta\theta\) ；由此推出 E 是 K 的次数为 \(q^e\) 的循环扩张，元素 \(\bar{\sigma}\) 生成 Gal(E/K)，并且 \(E = K(\theta)\) 。最后证明 K 的每个包含 F 的次数为 \(q^e\) 的循环扩张都是某个多项式 \(X^q - c\alpha\) （在 F{[}X{]} 中的，取适当的 \(c \in K^*\) ）的根域（利用 V 中第 85 页的定理 3）。
\item
  令 K 为有理数域 Q。多项式 \(X^{2} + 1\) 在 Q{[}X{]} 中不可约；如果 i 是它的一个根，则 \(F = \mathbf{Q}(i)\) 是 Q 的 2 次循环扩张，但不存在包含 F 的 Q 的 4 次循环扩张。
\end{enumerate}

\begin{enumerate}
\def\labelenumi{\arabic{enumi})}
\setcounter{enumi}{4}
\item
  设 K 是特征为 p 的域，n 是不为 p 的倍数的整数。对于 K{[}X{]} 中形如 \(X^{n} - a\) 的一个多项式，要使其不可约，必要且充分条件是：对于 n 的每个素因子 q，a 不等于 K 中某个元素的 q 次幂，并且当 \(n \equiv 0 \pmod{4}\) 时，a 不应具有形式 \(-4c^{4}\) （其中 \(c \in K\) ）。（为证明条件的充分性，化归到 \(n = q^{e}\) （q 为素数）的情形，并利用 V 章第 153 页习题 4。对 e 进行归纳论证，借助 V 章第 153 页习题 4 确定 \(X^{q^{e}} - a\) 的每个不可约因子的常数项的形式。）
\item
  设 \(\mathbf{K}\) 是特征为 \(p, n\) 一个不是 \(p\) 的倍数的整数，且 \(\mathbf{N}\) 为多项式 \(\mathbf{X}^n - a\) （ \(\mathbf{K}[\mathbf{X}]\) 的）的一个分裂域.
\end{enumerate}

\begin{enumerate}
\def\labelenumi{\alph{enumi})}
\tightlist
\item
  证明群 \(\operatorname{Gal}(\mathbf{N} / \mathbf{K})\) 同构于群 \(\Gamma\) （由矩阵 \(\left( \begin{array}{cc}x & y\\ 0 & 1 \end{array} \right)\) ，其中 \(x\in (\mathbf{Z} / n\mathbf{Z})^*\) 且 \(y\in \mathbf{Z} / n\mathbf{Z}\) ，所组成的）的一个子群。若 \(n\) 是素数，群
\end{enumerate}

Gal(N/K) 仅当 \(\mathrm{R}_{n}(\mathrm{K})=\mathrm{K}\) 或者 a 是 K 中某个元素的 n 次幂时才是交换的。b) 如果 K=Q 且 n 是素数，证明若 a 不是某个有理数的 n 次幂，则 Gal(N/K) 同构于整个 \(\Gamma\) .

\begin{enumerate}
\def\labelenumi{\alph{enumi})}
\setcounter{enumi}{2}
\tightlist
\item
  如果 \(\mathbf{K} = \mathbf{Q}\) ，确定 \(\operatorname{Gal}(\mathbf{N} / \mathbf{K})\) （当 \(a\) 是一个无平方因子整数时）。
\end{enumerate}

\begin{enumerate}
\def\labelenumi{\arabic{enumi})}
\setcounter{enumi}{6}
\tightlist
\item
  域 K 的一个扩张 E 称为库默尔扩张，如果存在整数 n \textgreater{} 0 使得 K 包含一个本原 n 次单位根，且 E 是伽罗瓦群被 n 零化的阿贝尔扩张。
\end{enumerate}

\begin{enumerate}
\def\labelenumi{\alph{enumi})}
\item
  设 K 为一个域，N 为 K 的 m 次伽罗瓦扩张，其伽罗瓦群为 \(\Gamma\) 。要使 N 成为 K 的库默尔扩张，必要且充分条件是存在 N 在 K 上的一个基 \((\theta_{\sigma})_{\sigma\in\Gamma}\) ，使得元素 \(\gamma_{\sigma\tau}=\theta_{\sigma}^{1-\tau}\) 对每一对 \((\sigma,\tau)\) 都属于 K。（对于条件的必要性，可归结为 N 是次数为素数幂的循环扩张的情形。对于充分性，注意若 \(\tau\in\Gamma\) 的阶为 n，则存在 \(\sigma\in\Gamma\) 使得 \(\gamma_{\sigma\tau}\) 是本原 n 次单位根。）
\item
  如果 \((\theta_{\sigma})_{\sigma \in \Gamma}\) 是满足条件 \(a)\) 的基，证明对于一个元素 \(x = \sum_{\sigma} a_{\sigma} \theta_{\sigma}\) （ \(N\) 的）
\end{enumerate}

使得 \(x\) 的共轭元构成 \(\mathbf{N}\) 在 \(\mathbf{K}\) 上的一个正规基，其充要条件是 \(a_{\sigma} \neq 0\) 对所有 \(\sigma \in \Gamma\) 成立.

* 8) 设 N 是域 K 的一个库默尔扩张。

\begin{enumerate}
\def\labelenumi{\alph{enumi})}
\item
  设 E 是 N 的一个子扩张，使得 \([E:K]=q\) 是一个素数。证明存在将 N 分解为 K 的循环扩张的分解 \(N=N_{1}\otimes N_{2}\otimes\cdots\otimes N_{r}\) ，使得 \(E\subset N_{j}\) 对某个索引 j 成立。
\item
  由 a) 可推出，如果 \(x \in N\) 使得 x 在 K 上的共轭元素构成 N 在 K 上的正规基，则对于每个满足 \(K \subset E \subset N\) 的域 E，x 在 E 上的共轭元构成 N 在 E 上的一个正规基（归结为 \([E : K]\) 是素数的情形，并利用习题 7，b）。*
\end{enumerate}

\begin{enumerate}
\def\labelenumi{\arabic{enumi})}
\setcounter{enumi}{8}
\tightlist
\item
  设 \(\mathbf{K}\) 是特征为 \(p > 0\) 的域.
\end{enumerate}

\begin{enumerate}
\def\labelenumi{\alph{enumi})}
\item
  设 E 是 K 的循环扩张，次数为 \(p^e\) ，并设 \(\sigma\) 是生成群 \(\operatorname{Gal}(\mathrm{E} / \mathrm{K})\) 的 K-自同构。设 F 是 K 上次数为 \(p^{e - 1}\) 的域，使得 \(\mathbf{K} \subset \mathbf{F} \subset \mathbf{E}\) ；如果 \(m = p^{e - 1}\) ，证明存在 \(\theta \in \mathbf{E}\) ，它是不可约多项式 \(\mathbf{X}^p - \mathbf{X} - \alpha\) （ \(\mathrm{F}[\mathbf{X}]\) 中的）的一个根，使得 \(\mathrm{E} = \mathrm{F}(\theta)\) 且 \(\sigma^m (\theta) = \theta + 1\) 。证明我们还有 \(\mathrm{E} = \mathrm{K}(\theta)\) 且 \(\sigma (\theta) = \theta + \beta\) ，其中 \(\beta \in \mathrm{F}\) 使得 \(\sigma (\alpha) - \alpha = \beta^p - \beta\) 且 \(\mathrm{Tr}_{\mathrm{F} / \mathrm{K}}(\beta) = 1\) .
\item
  反之，设 F 是 K 的次数为 \(p^{e-1}\) 的循环扩张（其中 e \textgreater{} 1), \(\sigma\) 是 F 的一个生成 Gal(F/K) 的 K-自同构。证明存在 F 的两个元素 \(\alpha\) , \(\beta\) ，使得 \(\operatorname{Tr}_{\mathrm{F}/\mathrm{K}}(\beta) = 1\) 且 \(\sigma(\alpha) - \alpha = \beta^{p} - \beta\) （利用 V 第 85 页定理 3 及 V 第 49 页命题 1 的推论）。由此推出对所有 \(c \in K\) ，多项式 \(X^{p} - X - \alpha - c\) 在 F{[}X{]} 中不可约；如果 \(\theta\) 是该多项式在 F 的一个代数闭包 \(\Omega\) 中的一个根，证明存在一个 K-同态 \(\bar{\sigma}\) ，它把 E = F( \(\theta\) ) 映到 \(\Omega\) 内，延拓 \(\sigma\) 并且满足 \(\bar{\sigma}(\theta) = \theta + \beta\) ；由此得出 E 是 K 的一个次数为 \(p^{e}\) 的循环扩张， \(\bar{\sigma}\) 生成 Gal(E/K) 且 E = K( \(\theta\) )。最后证明，K 的每个包含 F 的 \(p^{e}\) 次循环扩张都是多项式 \(X^{p} - X - \alpha - c\) （F{[}X{]} 中的，取适当的 \(c \in K\) ）的分裂域.
\end{enumerate}

\begin{enumerate}
\def\labelenumi{\arabic{enumi})}
\setcounter{enumi}{9}
\tightlist
\item
  \begin{enumerate}
  \def\labelenumii{\alph{enumii})}
  \tightlist
  \item
    设 \(\mathbf{K}\) 是一个域，其代数闭包是素数次数 \(q\) 的 \(\mathbf{K}\) 的扩张。这样 \(\mathbf{K}\) 是完全域，并且 \(q\) 必须不同于 \(\mathbf{K}\) 的特征（使用习题 9）。
  \end{enumerate}
\end{enumerate}

\begin{enumerate}
\def\labelenumi{\alph{enumi})}
\setcounter{enumi}{1}
\tightlist
\item
  证明 K 包含 q 次单位根，并且 \(\Omega\) 是不可约多项式 \(X^{q}-a\) （K{[}X{]} 中的）的分裂域；推断出 q=2（在相反情形下，由习题 5 推断该多项式 \(X^{q^{2}}-a\) 将是不可约的）。进一步证明 -a 在 K 中是平方（见上述引文），-1 在 K 中不是平方，并由此推断出 \(\Omega=\mathrm{K}(i)\) ，其中 \(i^{2}=-1\) .
\end{enumerate}

\begin{enumerate}
\def\labelenumi{\arabic{enumi})}
\setcounter{enumi}{10}
\item
  设 \(\mathbf{K}\) 是一个域，其代数闭包 \(\Omega\) 是一个有限次 \(> 1\) 的 \(\mathbf{K}\) 的扩张。证明 \(\Omega = \mathbf{K}(i)\) ，其中 \(i^2 = -1\) 。（如果我们有 \(\Omega \neq \mathbf{K}(i)\) ，借助伽罗瓦理论和西罗定理证明，那么就会存在一个域 \(E\) 使得 \(\mathbf{K}(i) \subset E \subset \Omega\) 且 \(\Omega\) 是 \(E\) 的素数次扩张；现在应用习题 10。）
\item
  设 K 是一个特征为 q 的域， \(\Omega\) 为 K 的一个代数闭扩张，x 为 \(\Omega\) 的一个元素，使得 \(x \notin K\) ，且 M 为 \(\Omega\) 的不包含 x 的子扩张所成集合的一个极大元（V，第 156 页，习题 9）；进一步假设 M 是完全域，使得 \(\mathbf{M}(x)\) 是素数次数 p 的 M 的伽罗瓦扩张（同上）。
\end{enumerate}

\begin{enumerate}
\def\labelenumi{\alph{enumi})}
\item
  证明：要么 \(p = 2\) 且 \(\Omega = M(i)\) ，要么对于每个整数 \(r \geqslant 1\) ，存在唯一一个次数为 \(p^r\) 的、关于 \(M\) 的扩张；该扩张 \(M_r\) 在 \(M\) 上是循环的， \(\Omega\) 是诸 \(M_r\) 的并，且诸 \(M_r\) 是 \(M\) 的仅有的有限次扩张。（注意，如果一个 \(p\) -群 \(\Gamma\) 的阶为 \(p^r\) ，且它只有一个阶为 \(p^{r-1}\) 的子群，则它必然是循环的；可对 \(r\) 作归纳论证，并利用 I，第 133 页，习题 10，以及 I，第 77 页，命题 11 的推论）。
\item
  假设 \(\Omega \neq \mathbf{M}(i)\) 且 \(p \neq q\) 。证明 \(\mathbf{M}\) 包含 \(p\) 次单位根；于是我们有 \(\mathbf{M}(x) = \mathbf{M}_1 = \mathbf{M}(\alpha)\) ，其中 \(\alpha^p \in \mathbf{M}\) 。证明若 \(\alpha\) 是一个 \(p\) 次幂（ \(\mathbf{M}_1\) 中某个元素的），则 \(\mathbf{M}_1 = \mathbf{M}(\zeta)\) ，其中 \(\zeta\) 是一个 \(p^2\) 次单位根（若 \(\alpha = \beta^p\) ，其中 \(\beta \in \mathbf{M}_1\) ，考虑 \(\beta\) 在 \(\mathbf{M}[\mathbf{X}]\) 中的极小多项式）。
\item
  在 \(b\) ) 的假设下，证明若 \(\alpha\) 不是一个 \(p\) 次幂（ \(\mathbf{M}_1\) 中某个元素的），则 \(\mathbf{M}_r\) 是多项式 \(\mathbf{X}^{p^{r-1}} - \alpha\) （ \(\mathbf{M}_1[\mathbf{X}]\) 中的）的分裂域（利用 \(b\) )).
\end{enumerate}

\(d)\) 在 \(b\) ) 的假设下，证明若 \(\alpha\) 是一个 \(p\) 次幂（ \(\mathbf{M}_1\) 中某个元素的），则存在 \(\gamma \in \mathbf{M}_1\) 使得 \(\gamma \notin \mathbf{M}\) 且 \(\mathbf{M}_2\) 是 \(\mathbf{X}^p - \gamma\) 的分裂域；证明 \(\gamma\) 不是一个 \(p\) 次幂（ \(\mathbf{M}_2\) 中某个元素的），并由此推出 \(\mathbf{M}_r\) 是多项式 \(\mathbf{X}^{p^{r-1}} - \gamma\) （ \(\mathbf{M}_1[\mathbf{X}]\) 中的）的分裂域。

\begin{enumerate}
\def\labelenumi{\arabic{enumi})}
\setcounter{enumi}{12}
\tightlist
\item
  如果域 K 的每一个在 K 上的有限次可分代数扩张都是循环扩张，则称 K 为 Moriya 域。* 每个有限域都是 Moriya 域（V，第 95 页，命题 4）。* 习题 12 中的域 M 是 Moriya 域。
\end{enumerate}

\begin{enumerate}
\def\labelenumi{\alph{enumi})}
\item
  要使 K 成为 Moriya 域，当且仅当对每个整数 n \textgreater{} 0，在 K 上至多存在一个 n 次可分扩张。（使用如下事实：若 G 是有限群，且对每个整除 G 的阶的整数 \(m \geqslant 1\) ，G 中至多存在一个 m 阶子群，则 G 是循环群；这可由 p-群的同一性质（习题 12）以及 I，第 80 页，定理 4 推出。）
\item
  若 K 是 Moriya 域，证明 K 的每个代数扩张 E 都是 Moriya 域，并且若 F 是 E 上的一个 m 次扩张，则存在 K 上的一个 m 次扩张。
\end{enumerate}

\begin{enumerate}
\def\labelenumi{\arabic{enumi})}
\setcounter{enumi}{13}
\tightlist
\item
  \begin{enumerate}
  \def\labelenumii{\alph{enumii})}
  \tightlist
  \item
    设 K 是一个域，使得存在 K 的任意大有限次的代数扩张，且所有这些扩张的次数都是同一个素数 p 的幂。证明在这些条件下，对每个幂 \(p^{h}\) ( \(h \geqslant 0\) ）都存在 K 的一个次数为 \(p^{h}\) 的扩张（利用 I，第 77 页，定理 1）。
  \end{enumerate}
\end{enumerate}

\begin{enumerate}
\def\labelenumi{\alph{enumi})}
\setcounter{enumi}{1}
\tightlist
\item
  设 p 为一个素数，K 为一个域，使得存在 K 的一个代数扩张，其次数在 \(p \neq 2\) 时被 p 整除，在 p = 2 时被 4 整除。证明存在 K 的一个代数扩张 E，使得 E 的代数扩张的次数的集合恰为 p 的幂 \(p^{h}\) 的集合。（若 K 是完全域或特征 \(\neq p\) ，令 \(K_{1}\) 为 K 的完全闭包；考虑 \(K_{1}\) 的一个极大的代数扩张，在所有其元素在 \(K_{1}\) 上的次数不被 p 整除的那些代数扩张之中选取；利用 a) 和习题 11。）
\end{enumerate}

\begin{enumerate}
\def\labelenumi{\arabic{enumi})}
\setcounter{enumi}{14}
\item
  设 K 为 Moriya 域（习题 13），P 为整除 K 的有限次代数扩张的次数的素数集合。证明：K 的有限次代数扩张的次数的集合，要么等于集合 \(N_{P}\) （即其所有素因子均在 P 中的整数所构成的集合），要么等于 \(N_{P}\) 的由不能被 4 整除的整数组成的子集。
\item
  对于每个素数 \(p\) ，存在一个特征为 \(p\) 的 Moriya 域，使得该域的有限次代数扩张的次数的集合是集合 \(\mathbf{N}^*\) （即所有 \(>0^1\) 的整数所组成的集合）。证明对于素数的每一个集合 \(P\) ，存在一个任意特征的 Moriya 域 \(K\) ，使得 \(K\) 的有限次代数扩张的次数的集合是集合 \(N_P\) （如习题 14, b) 中那样论证）。
\end{enumerate}

* 17) 设 \(\mathbf{P}_n\) 为首一多项式 \(\mathbf{X}^n + a_1\mathbf{X}^{n-1} + \cdots + a_n\) （ \(\mathbf{Z}[\mathbf{X}]\) 中的）所成的集合，这些多项式在 \(\mathbf{C}\) 中的所有根都具有绝对值 \(\leqslant 1\) .

\begin{enumerate}
\def\labelenumi{\alph{enumi})}
\item
  证明集合 \(P_{n}\) 是有限的。由此推出，对于每个多项式 \(F \in P_{n}\) ，如果 \(F_{h}\) 表示以F的根的h次幂为根的多项式，则存在两个不同的整数h、k，使得 \(F_{h} = F_{k}\) .
\item
  由 a) 推出，多项式 \(F \in P_{n}\) 的所有根要么为零，要么是单位根（“克罗内克定理”）。*
\end{enumerate}

\begin{enumerate}
\def\labelenumi{\arabic{enumi})}
\setcounter{enumi}{17}
\tightlist
\item
  设 \(p \in \mathbb{N}\) 是一个素数， \(C_p\) 为级为 \(p\) 的、关于 \(\mathbf{Q}\) 的分圆扩张， \(\zeta \in C_p\) 为一个本原 \(p\) 次单位根， \(n\) 为一个整数。记
\end{enumerate}

\[
\mathrm{X} _ {n} (\zeta) = \sum_ {\alpha \in \mathbf {Z} / p \mathbf {Z}} \zeta^ {\alpha^ {n}}.
\]

\begin{enumerate}
\def\labelenumi{\alph{enumi})}
\item
  计算 \(d = [\mathbf{Q}(\mathbf{X}_n(\zeta)) : \mathbf{Q}]\) 。确定扩张 \(\mathbf{Q}(\mathbf{X}_n(\zeta))\) （ \(\mathbf{Q}\) 上的）的伽罗瓦群，以及 \(\mathbf{X}_n(\zeta)\) 的共轭元。
\item
  设 \(T^{d} + a_{1}T^{d-1} + \cdots + a_{d}\) 为 \(\mathbf{X}_{n}(\zeta)\) 在 Q 上的极小多项式。证明 \(a_{1} = 0\) .
\item
  假设 \(n > 1\) 且 \(p \equiv 1 \pmod{n}\) 。证明 \(a_2 = -\frac{n(n - 1)p}{2}\) （当 \(\frac{p - 1}{n}\) 为偶数时），而 \(a_2 = \frac{np}{2}\) （当 \(\frac{p - 1}{n}\) 为奇数时）。
\item
  假设 \(p \equiv 1 \pmod{n}\) 。设 \(\beta \in \mathbf{F}_p^*\) 是这样一个元素，其在 \(\mathbf{F}_p^* / \mathbf{F}_p^{*n}\) 中的剩余类是一个生成元。设 \(\mu\) 为方程在 \(\mathbf{F}_p^n\) 中的解的个数，该方程为 \(\sum_{i=1}^{n} \beta^{i-1} \mathbf{X}_i^n = 0\) 。证明 \(a_n = \frac{(-1)^n p(\mu - p^{n-1})}{p-1}\) .
\item
  计算 \(a_3\) ，其中 \(n = 3\) 且 \(p = 7, 13, 19\) .
\end{enumerate}

*1 域 \(\mathbf{Q}_p\) （即 \(p\) -进数域）在其代数闭包中的最大完全分歧扩张是一个具有相同性质的特征为 0 的域。*

* 19) a) 假设素数定理成立；如果 \(\pi(x)\) 是满足 \(p \leqslant x\) 的素数的个数，则 \(\pi(x) \sim x / \log x\) （当 \(x\) 趋于 \(+\infty\) ）。由此推出，对于每个奇数 \(m\) ，存在一个递增的素数序列 \(p_1 < p_2 < \cdots < p_m\) 使得 \(p_m < p_1 + p_2\) （考虑介于 \(x\) 与 \(3x/2\) 之间的素数）。

\begin{enumerate}
\def\labelenumi{\alph{enumi})}
\setcounter{enumi}{1}
\tightlist
\item
  利用由 a) 确定的有限素数序列 \(p_{j}\) ( \(1 \leq j \leq m\) ） ，令 \(n = p_{1}p_{2} \ldots p_{m}\) 。证明在分圆多项式 \(\Phi_{n}(X)\) 中，次数为 \(< p_{m} + 1\) 的项与多项式
\end{enumerate}

\[
(1 - \mathbf {X}) ^ {- 1} \prod_ {j = 1} ^ {m} (1 - \mathbf {X} ^ {p _ {j}}).
\]

的相应项相同。由此推出 \(\mathbf{X}^{p_m}\) 在 \(\Phi_n(\mathbf{X})\) 中的系数是 \(1 - m\) （“I. Schur 定理”）。*

* 20) 设 \(p\) 是一个素数，且 \(G\) 为一个交换 \(p\) -挠群（VII，第 10 页）。\\
a) 设 \(G_i \subset G\) 为满足以下条件的元素 \(x \in G\) 的集合：对每个 \(l\) ，方程 \(p^{l} y = x\) 有一个解。证明 \(G_i\) 是 \(G\) 的一个直因子。

\begin{enumerate}
\def\labelenumi{\alph{enumi})}
\setcounter{enumi}{1}
\item
  假设方程 px = 0 在 G 中只有有限多个解。证明：对每个整数 l，核 \(_{l}G\) （即 G 中乘以 \(p^{l}\) 的核）是一个有限群，并且秩 \(r(l)\) （即 \(F_{p}\) -向量空间 \(_{l}G/_{(l-1)}G\) 的秩）是 l 的递减函数，因此当 \(l > l_{0}\) （ \(l_{0}\) 足够大）时为常数。（利用初等因子定理（VII，第 24 页）。证明 \(G/G_{i}\) 是一个被 \(p^{l_{0}}\) 零化的有限群。）
\item
  进一步假设 \(\mathbf{G} = \mathbf{G}_i\) 。证明映射 \(l \mapsto r(l)\) 是常数。记 \(\mathbf{H} = (\mathbf{Z}[1/p]/\mathbf{Z})^{(1)}\) ，证明对每个 \(l\) 存在同构 \(_l\mathbf{G} \to {}_l\mathbf{H}\) ，并且每个同构 \(_l\mathbf{G} \to {}_l\mathbf{H}\) 扩展为同构 \((l+1)\mathbf{G} \to {}_{(l+1)}\mathbf{H}\) 。由此推出 \(\mathbf{G}\) 同构于 \(\mathbf{H}\) .
\item
  在 b) 的假设下，证明存在整数 n 和一个映射 \(i \to n_i\) （从 \(N - \{0\}\) 到 N 中，具有有限支撑），使得 G 同构于 \((\mathbf{Z}[1/p]/\mathbf{Z})^n \oplus \bigoplus_{i} (\mathbf{Z}/p^i \mathbf{Z})^{n_i}\) 。证明整数 n 和映射 \(i \mapsto n_i\) 是唯一地
\end{enumerate}

由映射 \(l \mapsto \text{Card}(l\mathbf{G})\) 所决定的。*

* 21) 设 G 和 H 是两个交换挠群，使得对每个整数 \(n\) ，方程 \(nx = 0\) 在 G 和 H 中具有相同的有限个解。证明 G 与 H 同构。（将 G 和 H 分解为 \(p\) -挠群的直和并应用习题 20。）令 \(G = \bigoplus_{i \in \mathbf{N}} \mathbf{Z} / p^i\mathbf{Z}\) 且 \(H = \bigoplus_{i \in \mathbf{N}} \mathbf{Z} / p^{2i}\mathbf{Z}\) 。证明 G 与 H 不同构

并且对每个整数 \(n\) ，方程 \(nx = 0\) 的解集的基数在 \(G\) 和 \(H\) 中是相同的。*

\begin{enumerate}
\def\labelenumi{\arabic{enumi})}
\setcounter{enumi}{21}
\tightlist
\item
  \begin{enumerate}
  \def\labelenumii{\alph{enumii})}
  \tightlist
  \item
    设 k 是一个域，f 是 \(k[X]\) 中的一个可分的首一多项式，D 是 f 的一个分裂扩张， \(f, x_{1}, \ldots, x_{n}\) 是 f 在 D 中的根，且 G 是 D 在 k 上的伽罗瓦群。假设 G 是循环群。证明：要使 G 的每个元素都诱导 \(x_{i}\) 的一个偶置换，其必要且充分条件是 f 的不可约因子的个数与 f 的次数模 2 同余（考虑 G 的一个生成元 \(\sigma\) 以及 \(x_{i}\) 的置换（由 \(\sigma\) 所诱导的）的轮换分解）。
  \end{enumerate}
\end{enumerate}

\begin{enumerate}
\def\labelenumi{\alph{enumi})}
\setcounter{enumi}{1}
\tightlist
\item
  进一步假设 k 的特征为 \(\neq 2\) 。证明上述两个条件等价于 f 的判别式是 k 中的平方这一事实（“Stickelberger 定理”）。
\end{enumerate}

\begin{enumerate}
\def\labelenumi{\arabic{enumi})}
\setcounter{enumi}{22}
\tightlist
\item
  设 \(n\) 和 \(p\) 是两个奇且互异的素数。
\end{enumerate}

\begin{enumerate}
\def\labelenumi{\alph{enumi})}
\item
  设 \(Q \in F_{p}[X]\) 为分圆多项式 \(\Phi_{n}\) 的像。证明：Q 的不可约因子的个数 i 为偶数当且仅当 p 是模 n 的平方。（设 K 为 Q 的分裂域。证明 \(\operatorname{Gal}(K/F_{p})\) 可等同于 \((\mathbf{F}_{n})^{*}\) 的由 p 的像生成的子群。研究指标 \(((\mathbf{F}_{n})^{*}: \operatorname{Gal}(K/F_{p}))\) 的奇偶性。）
\item
  证明 \(\mathbf{X}^n - 1 \in \mathbf{F}_p[\mathbf{X}]\) 的判别式是
\end{enumerate}

\[
(- 1) ^ {\frac {n (n - 1)}{2} + n - 1} n ^ {n}.
\]

\begin{enumerate}
\def\labelenumi{\alph{enumi})}
\setcounter{enumi}{2}
\tightlist
\item
  证明： \(i\) 是偶数当且仅当 \((-1)^{\frac{n(n - 1)}{2}} n\) 是模 \(p\) 的平方（把习题 22 应用于多项式 \(\mathbf{X}^n - 1\) ）。
\end{enumerate}

\[
\left(\frac {n}{p}\right) = 1
\]

\[
\left(\frac {n}{p}\right) = - 1
\]

\(\left(\frac{n}{p}\right)\left(\frac{p}{n}\right) = (-1)^{\frac{(n - 1)(p - 1)}{4}}\) （“二次互反律”）。
